\documentclass[11pt]{article}

\usepackage{acl}
\usepackage{times}
\usepackage{latexsym}
\usepackage{microtype}
\usepackage{inconsolata}
\usepackage{amsmath}
\usepackage{amssymb}
\usepackage{mathtools}
\usepackage{booktabs}
\usepackage{graphicx}
\usepackage{subcaption}
\usepackage{multirow}
\usepackage{xcolor}
\usepackage{algorithm}
\usepackage{algpseudocode}
\usepackage{hyperref}
\usepackage{url}
\usepackage{bm}
\usepackage{enumitem}
\usepackage{wrapfig}
\usepackage{colortbl}
\usepackage{array}
\usepackage{makecell}
\usepackage{tikz}
\usetikzlibrary{shapes.geometric,arrows.meta,positioning,fit,backgrounds,calc}

\hypersetup{
  colorlinks=true,
  linkcolor=blue,
  citecolor=blue,
  urlcolor=blue
}
% ─── Float placement tuning ───────────────────────────────────────────────────
\renewcommand{\topfraction}{0.9}
\renewcommand{\bottomfraction}{0.8}
\renewcommand{\textfraction}{0.1}
\renewcommand{\floatpagefraction}{0.85}
\setcounter{topnumber}{4}
\setcounter{bottomnumber}{4}
\setcounter{totalnumber}{6}

% ─── Convenience macros ────────────────────────────────────────────────────────

\newcommand{\As}{A_{\mathrm{s}}}
\newcommand{\Ag}{A_{\mathrm{g}}}
\newcommand{\Rk}{R_k}
\newcommand{\mk}{m_k}
\newcommand{\xt}{x_t}
\newcommand{\zt}{z_t}
\newcommand{\zctrl}{z^{\mathrm{ctrl}}_{t}}
\newcommand{\zedit}{z^{\mathrm{edit}}_{R}}
\newcommand{\Itat}{I_{t}}
\newcommand{\pith}{\pi_\theta}
\newcommand{\tauP}{\tau_P}
\newcommand{\tauF}{\tau_F}
\newcommand{\R}{\mathbb{R}}
\newcommand{\eg}{\textit{e.g.}}
\newcommand{\ie}{\textit{i.e.}}
\newcommand{\etal}{\textit{et al.}}
\newcommand{\guardpaint}{GuardPaint}


\definecolor{passgreen}{HTML}{2ecc71}
\definecolor{failred}{HTML}{e74c3c}
\definecolor{lightgray}{HTML}{f2f2f2}




% Preamble
\usepackage{booktabs}
\usepackage{graphicx}
\usepackage{xcolor}
\usepackage{pifont}
\newcommand{\cmark}{\ding{51}}
\newcommand{\xmark}{\ding{55}}

% ─── Title & Authors ───────────────────────────────────────────────────────────
\title{
%\tipai{}--\tspo{}: Speculative Safety Decoding for\\
%Text-to-Image Generation via Tournament Inpainting
GuardPaint: Speculative Safety Decoding for Text-to-Image Generation\thanks{To further research in the field we release our code  \href{https://anonymous.4open.science/r/TiPAI-TSPO-D186/README.md}{here}.$^{\dagger}$This work was conducted outside the authors’ primary
professional roles.}
}


\author{
\textbf{
Shreyash Dhoot$^{1}$,
Paras Dhiman$^{1}$,
Arsh Abbas Naqvi$^{1}$,
Arnabi Dutta$^{1}$
}\\
\textbf{
Aman Chadha$^{2,\dagger}$,
Vinija Jain$^{3,\dagger}$,
Amitava Das$^{1}$
}\\[4pt]
\normalfont
$^{2}$Apple, USA \quad
$^{3}$Meta, USA  \quad
$^{1}$Pragya Lab, BITS Pilani Goa, India 
}


\begin{document}
\maketitle

%% ════════════════════════════════════════════════════════════════════════════
\begin{abstract}
%% ════════════════════════════════════════════════════════════════════════════

Text-to-image (T2I) diffusion models offer powerful visual generation, but their controllability creates a critical safety challenge: adversarial prompts can steer the denoising trajectory toward policy-violating content such as explicit nudity or graphic violence. Existing safeguards mostly act before generation through prompt filtering or after generation through image classification, leaving the diffusion process itself unguarded and often yielding only refusal rather than safe visual repair.

We introduce GuardPaint, a speculative decoding framework for safe T2I generation that intervenes inside the diffusion trajectory without modifying the base model. A lightweight auditor monitors intermediate images, localizes unsafe regions, and triggers surgical inpainting repair only where needed. Candidate repairs are generated by a policy-aligned inpainter and selected through a guarded tournament that accepts edits only when they improve policy compliance while preserving prompt fidelity and perceptual quality.

Across five jailbreak families SneakPrompt, MMA, PGJ, DACA, and RABell and UNet/flow-matching models including SD~1.5, SDXL, SD~3.5, and FLUX.1-dev. GuardPaint reduces attack success and harmful generations with minimal degradation to image quality, prompt fidelity, and benign behavior.

Content warning: This paper contains examples involving nudity and violence that some readers may find disturbing, distressing, or offensive.

\end{abstract}









\begin{table*}[ht!]
\centering
\small
\setlength{\tabcolsep}{3.8pt}
\renewcommand{\arraystretch}{1.18}
\caption{\textbf{\emph{Positioning of \textsc{GuardPaint}.}}
Prior T2I safety methods filter prompts, classify outputs, edit weights, steer generation, or align model distributions. \textsc{GuardPaint} combines trajectory-level intervention, localized repair, safe visual alternatives, plug-and-play deployment, and tournament-based repair selection.}
\label{tab:guardpaint_positioning}

\begin{tabular}{p{0.43\textwidth}ccccc}
\toprule
\textbf{Method family} &
\textbf{\makecell{Trajectory\\level}} &
\textbf{\makecell{Localized\\repair}} &
\textbf{\makecell{Safe visual\\alternative}} &
\textbf{\makecell{No base\\retraining}} &
\textbf{\makecell{Tournament\\selection}} \\
\midrule
Prompt filtering / rewriting~\cite{posi2024,lyu2024latentguard} &
\xmark & \xmark & \xmark & \cmark & \xmark \\
Post-hoc safety classifiers &
\xmark & \xmark & \xmark & \cmark & \xmark \\
Concept erasure / model editing~\cite{schramowski2023safe,gandikota2023erasing} &
\xmark & \xmark & \xmark & \xmark & \xmark \\
Inference-time safety steering~\cite{meng2026safetytaxmitigatingunsafe,kim2025safedenoiser} &
\cmark & \xmark & \cmark & \cmark & \xmark \\
Preference-aligned diffusion models~\cite{wallace2023diffusion_dpo,li2024diffusion_kto,borso2025d3po,hong2025marginawarepreferenceoptimizationaligning,duo2024,li2024safegen,zhang2025shielddiff} &
\xmark & \xmark & \cmark & \xmark & \xmark \\
Preference-tuned inpainting~\cite{liu2024prefpaint} &
\xmark & \cmark & \cmark & \cmark & \xmark \\
\midrule
\textbf{\textsc{GuardPaint}} &
\cmark & \cmark & \cmark & \cmark & \cmark \\
\bottomrule
\end{tabular}
\end{table*}

%% ════════════════════════════════════════════════════════════════════════════
%% %% ════════════════════════════════════════════════════════════════════════════
\section{The Case for Plug-and-Play Safety Alignment in Text-to-Image Generation}
\label{sec:motivation_positioning}
%% ════════════════════════════════════════════════════════════════════════════

Text-to-image (T2I) diffusion models now produce high-fidelity images from open-ended prompts, but the same controllability creates a sharp safety failure mode: \textbf{\emph{adversarial prompts can steer the denoising trajectory toward policy-violating content}}, including explicit nudity and graphic violence. Recent attacks exploit prompt obfuscation, token substitution, multilingual or Unicode perturbations, and semantic rewriting to bypass surface-level safeguards~\cite{yang2023sneakprompt,mma2024,rabell2024,daca2024,pgj2024,struppek2023homoglyphs,ba2023surrogateprompt}. As T2I systems move into creative and public-facing deployment, safety mechanisms must block harmful content \textbf{\emph{without unnecessarily degrading benign generations}}.

Existing defenses mostly protect the \emph{boundaries} of generation. Prompt-level safeguards detect, rewrite, or normalize unsafe inputs before sampling~\cite{posi2024,lyu2024latentguard}, while post-hoc classifiers intervene only after an image has already been produced. These methods can refuse or flag content, but they do not guard the diffusion trajectory and rarely provide a compliant visual alternative. Other methods modify the generator itself through concept erasure, model editing, safe denoising, inference-time steering, or preference alignment~\cite{schramowski2023safe,gandikota2023erasing,meng2026safetytaxmitigatingunsafe,kim2025safedenoiser,wallace2023diffusion_dpo,li2024diffusion_kto,borso2025d3po,hong2025marginawarepreferenceoptimizationaligning,duo2024,li2024safegen,zhang2025shielddiff,liu2024prefpaint}. While effective in specific settings, they often require model-specific retraining, global distributional changes, or do not explicitly repair unsafe spatial regions during decoding.

This motivates a central systems question:
\begin{quote}
\centering
\emph{\textbf{Can safety alignment be introduced as a modular decoding-time layer without modifying the underlying generator?}}
\end{quote}

We argue that such a layer requires moving from \emph{boundary-level filtering} to \textbf{\emph{trajectory-level repair}}. Rather than rejecting prompts before sampling or classifying images after generation, a safety layer should monitor the denoising process itself, detect when unsafe structure emerges, and intervene locally. The desired intervention is \textbf{\emph{plug-and-play}} across evolving T2I backbones, \textbf{\emph{localized}} to unsafe regions, and \textbf{\emph{selective}} enough to reject edits that harm prompt fidelity or perceptual quality.

\textsc{GuardPaint} follows this premise. Inspired by speculative decoding in language models~\cite{leviathan2023speculative,chen2023accelerating}, the frozen diffusion generator acts as a draft model and a lightweight auditor acts as a verifier. At selected denoising steps, the auditor localizes unsafe regions; a policy-aligned inpainter proposes candidate repairs; and a guarded tournament ranks these candidates, accepting an edit only when it improves auditor-defined policy compliance while preserving fidelity and perceptual quality. The result is a \textbf{\emph{plug-and-play safety alignment layer}} that converts unsafe generations into safe, semantically coherent visual alternatives rather than blank refusals.

\paragraph{Contributions.}
\begin{itemize}[leftmargin=1.2em,itemsep=3pt,topsep=3pt]
    \item We introduce \textbf{\emph{plug-and-play safety alignment}} for T2I diffusion models via decoding-time trajectory intervention instead of generator retraining.
    \item We propose \textbf{\emph{\textsc{GuardPaint}}}, a speculative safety decoding framework combining adversarial auditing, localized inpainting repair, and guarded candidate ranking.
    \item We introduce a \textbf{\emph{guarded tournament alignment mechanism}} that accepts repairs only when they improve auditor-defined policy compliance while preserving fidelity and perceptual quality.
    \item We demonstrate deployment across UNet-based and flow-matching T2I architectures without modifying base-model weights.
\end{itemize}



%% %% ════════════════════════════════════════════════════════════════════════════
\section{GuardPaint: Speculative Safety Decoding}
\label{sec:framework}
%% ════════════════════════════════════════════════════════════════════════════

\begin{figure*}[t]
  \centering
  \includegraphics[width=\textwidth]{inference.png}
  \caption{
    \textbf{\textsc{GuardPaint} inference pipeline.}
    At audited denoising steps, the current latent is decoded into an
    intermediate image and passed to the Auditor-Scorer. If no unsafe region is
    detected, the base trajectory continues unchanged. If a suspect region is
    detected, the tournament policy proposes $N$ inpainting configurations, the
    policy-aligned inpainter generates $N$ local repair candidates, and the
    guarded tournament selects a winner for latent reinsertion.
  }
  \label{fig:pipeline}
\end{figure*}

\paragraph{Overview.}
Let $p$ be the prompt and let
$\{z_t\}_{t=T}^{0}$ denote the denoising trajectory of a frozen generator
$\Phi_{\mathrm{base}}$. \textsc{GuardPaint} treats this trajectory as a
sequence of \textbf{\emph{auditable intermediate states}}. At selected
timesteps $t\in\mathcal{T}_{\mathrm{audit}}$, the latent is decoded into an
audit view $I_t=\mathrm{Dec}(z_t)$. A lightweight Auditor-Scorer $\As$
estimates three quantities: whether intervention is needed, where unsafe
content is localized, and whether a proposed repair preserves safety,
faithfulness, and perceptual quality. If no violation is detected, denoising
continues unchanged. If a violation is detected, a policy-aligned inpainter
$\Ag$ proposes local repairs, and a guarded tournament accepts an edit only
when it improves the calibrated safety objective while passing fidelity and
quality gates.

This yields a \textbf{\emph{draft--audit--repair--verify}} loop. The frozen
diffusion model is the \emph{draft}; the auditor is the \emph{verifier}; the
inpainter proposes \emph{localized repairs}; and the tournament performs
\emph{selective acceptance}. The key invariant is an
\textbf{\emph{auditor-defined non-regression guard}}: when no candidate passes
the guard, the base trajectory is preserved.

% ─────────────────────────────────────────────────────────────────────────────
\subsection{Auditor-Scorer: Detecting When and Where to Intervene}
\label{sec:auditor}
% ─────────────────────────────────────────────────────────────────────────────

The Auditor-Scorer $\As$ is a multi-task model with three roles:
\textbf{\emph{trigger}} repair, \textbf{\emph{localize}} unsafe regions, and
\textbf{\emph{score}} candidate repairs. It uses a ResNet-101 backbone
pretrained on ImageNet~\cite{he2016deep}.

\paragraph{Prompt and timestep conditioning.}
A lightweight BiLSTM text encoder with word embeddings
($d_{\mathrm{text}}=512$) encodes the prompt. Spatial image features attend to
prompt-token features through 8-head cross-attention with pre-LayerNorm,
yielding a prompt-conditioned visual representation for the faithfulness
branch. A 3-layer timestep MLP,
$128\!\to\!256\!\to\!512$ with SiLU activations, embeds
$t/T\in[0,1]$. Two FiLM projections~\cite{perez2018film} produce
$(\gamma,\beta)$ modulation parameters, allowing the auditor to calibrate
risk differently for early noisy views and late semantically formed images.

\paragraph{Multi-task outputs.}
The auditor produces six outputs, organized by function:
\begin{enumerate}[leftmargin=1.5em,itemsep=2pt,topsep=2pt]
  \item \textbf{\emph{Triggering:}} binary adversarial probability
  $\hat y_{\mathrm{adv}}\in[0,1]$ and class logits over
  $\{\mathrm{safe},\mathrm{nudity},\mathrm{violence}\}$.
  \item \textbf{\emph{Localization:}} per-class risk maps
  $r_c\in[0,1]^{H\times W}$ used to mine unsafe spatial regions.
  \item \textbf{\emph{Tournament scoring:}} prompt faithfulness $F$,
  seam/perceptual quality $P$, and relative adversarial suppression strength
  $B\in[0,1]$.
\end{enumerate}
The faithfulness branch is isolated from safety-label gradients: safety labels
do not update the image--text alignment head, preserving a clean fidelity
signal for tournament gating.

\paragraph{Region mining.}
At inference, the adversarial map head
$\mathrm{Conv2d}(2048,1,1)\!\to\!\mathrm{Sigmoid}$ produces
$r_{\mathrm{adv}}\in[0,1]^{H'\times W'}$. This map is bilinearly upsampled to
$512\times512$, thresholded at the 85th percentile, and feathered with a
Gaussian kernel $(\sigma=5,\;15\times15)$ to smooth mask boundaries before
inpainting.

\paragraph{Training objective.}
The auditor is trained with a \textbf{\emph{factorized multi-task objective}}
that separates triggering, harm typing, repair quality, and prompt fidelity:
\begin{align*}
\mathcal{L}_{\As}
&=
\underbrace{\mathcal{L}_{\mathrm{adv}}}_{\text{repair trigger}}
+
0.5\underbrace{\mathcal{L}_{\mathrm{class}}}_{\text{harm type}}
+
0.4\underbrace{\mathcal{L}_{\mathrm{rel\text{-}adv}}}_{\text{suppression score}}
\\[-0.15em]
&\quad+
0.3\underbrace{\mathcal{L}_{\mathrm{seam}}}_{\text{visual quality}}
+
0.5\underbrace{\mathcal{L}_{\mathrm{InfoNCE}}}_{\text{prompt fidelity}} .
\end{align*}
Here $\mathcal{L}_{\mathrm{adv}}$ gates whether repair is triggered;
$\mathcal{L}_{\mathrm{class}}$ selects the harm category and risk map;
$\mathcal{L}_{\mathrm{rel\text{-}adv}}$ estimates adversarial strength;
$\mathcal{L}_{\mathrm{seam}}$ penalizes boundary artifacts; and
$\mathcal{L}_{\mathrm{InfoNCE}}$~\cite{oord2018cpc} supervises image--text
faithfulness independently of safety labels.

\paragraph{Dataset and labeling.}
We assemble 82,000 image--prompt pairs from eight public HuggingFace datasets
covering safe, nudity, and violence categories. Images without prompts are
captioned using Qwen2.5-VL~\cite{bai2025qwen25vltechnicalreport}. All images
are relabeled with InternVL-3.5 (8B)~\cite{internvl3p5} into
$\{\mathrm{safe},\mathrm{nudity},\mathrm{violence}\}$, selected after
qualitative labeler ablations against four alternative VLMs
(Table ~\ref{tab:labeler_ablation}). We use a 70/15/15 split; dataset
sources appear in Table~\ref{tab:dataset}.

\begin{table}[t]
\centering
\small
\setlength{\tabcolsep}{5pt}
\begin{tabular}{lccc}
\toprule
\textbf{Class} & \textbf{P} & \textbf{R} & \textbf{F1} \\
\midrule
Safe     & 0.90 & 0.97 & 0.94 \\
Nudity   & 0.90 & 0.78 & 0.84 \\
Violence & 0.85 & 0.89 & 0.87 \\
\midrule
Macro avg & 0.88 & 0.88 & 0.88 \\
\bottomrule
\end{tabular}
\caption{Auditor-Scorer classification performance on the held-out test set
(overall accuracy: 88\%).}
\label{tab:auditor_perf}
\end{table}

% ─────────────────────────────────────────────────────────────────────────────
\subsection{Policy-Aligned Inpainter: Learning Safe Local Repairs}
\label{sec:inpainter}
% ─────────────────────────────────────────────────────────────────────────────

The inpainter $\Ag$ converts a masked unsafe region into a visually coherent,
policy-compliant alternative. We train it in two stages. The first stage
teaches \textbf{\emph{what safe repair looks like}}; the second teaches
\textbf{\emph{when safe repair should dominate unsafe completion}}. This
mirrors SFT-then-alignment in language-model safety and follows the staged
preference-tuning pattern used in inpainting~\cite{liu2024prefpaint}.

\paragraph{Stage 1: Refusal supervised fine-tuning.}
We start from Stable Diffusion 1.5 Inpainting and apply LoRA
fine-tuning~\cite{hu2022lora} with $r=128$ and $\alpha=128$ on attention
modules, feed-forward layers, spatial projection layers, and the UNet input
convolution. Training examples pair unsafe prompts and masked unsafe images
with safe target completions, using masks derived from auditor heatmaps. A
Min-SNR weighted MSE loss~\cite{hang2023minsnr} is applied across
$t\in[0,1000]$, teaching the model a stable \textbf{\emph{safe repair
manifold}} for masked adversarial regions. The LoRA weights are merged into
the base checkpoint before alignment.

\paragraph{Stage 2: Binary Classifier Optimization.}
On the SFT-merged model, we apply Binary Classifier Optimization
(BCO)~\cite{jung2024bco} with binary safe/unsafe labels. Each training step
samples a clean latent $z_0$, timestep $t$, and Gaussian noise $\epsilon$:
\begin{align*}
t &\sim \mathcal{U}[0,1000),
\qquad
\epsilon\sim\mathcal{N}(0,I),                                      \\[-0.1em]
z_t
&=
\sqrt{\bar\alpha_t}\,z_0
+
\sqrt{1-\bar\alpha_t}\,\epsilon .
\end{align*}
The trainable denoiser $\epsilon_\theta$ and frozen reference
$\epsilon_{\mathrm{ref}}$ receive the same tuple
$(z_t,t,m_\ell,\tilde z_0,p)$, where $\tilde z_0$ is the masked latent and
$m_\ell$ is the latent-space mask.

\paragraph{Masked reconstruction reward.}
Both denoisers are converted into clean-latent estimates
$\hat z^\theta_0$ and $\hat z^{\mathrm{ref}}_0$. Alignment is driven by the
masked reconstruction gap:
\begin{align*}
g
&=
\ell_{\mathrm{ref}}-\ell_{\theta},\\
\ell_{\bullet}
&=
\frac{
\sum
\bigl(1+\lambda_m m_\ell\bigr)
\bigl(\hat z^{\bullet}_0-z_0\bigr)^2m_\ell}{\sum m_\ell}, 
\lambda_m=0.5 .
\end{align*}
A positive $g$ means the trainable model reconstructs the masked region better
than the reference. This is rewarded for safe examples and penalized for
unsafe ones.

\paragraph{BCO loss.}
An EMA baseline stabilizes the reward boundary:
\begin{align*}
r
&=
g-\delta_{\mathrm{ema}},
\qquad
\delta_{\mathrm{ema}}\in[-0.03,+0.03],                              \\[-0.1em]
\mathrm{sgn}
&=
2y-1,
\qquad
y\in\{0,1\}.
\end{align*}
The BCO objective is
\begin{align*}
\mathcal{L}_{\mathrm{BCO}}
&=
\mathbb{E}\!\left[
\mathrm{softplus}
\left(
-\mathrm{sgn}\cdot\beta r
\right)
\right],
\qquad
\beta=50 .
\end{align*}
We use class weights 1.0/5.0/12.0 for safe/nudity/violence, imposing higher
cost on unsafe misses. To reduce the single-step/multi-step mismatch, every
10 gradient steps we run $K{-}1$ deterministic DDIM steps~\cite{song2020ddim}
under \texttt{no\_grad}, re-noise to
$t_{\mathrm{final}}=\lfloor t/K\rfloor$, and apply a second BCO update at the
lower noise level. This teaches repair at the stage where unsafe structure is
visually coherent. Full details are in Appendix~\ref{app:inpainter_full}.

\begin{figure}[!ht]
  \centering
  \includegraphics[width=\linewidth]{final-final-dig.png}
  \caption{
    \textbf{Tournament example}: rows show examples of proposed inpainted candidates in a tournament for different adversarial prompts. The yellow border marks the \emph{WINNER} candidate for that tournament round.(Translucent red mask shows the the binary mask/the inpainted area)
  }
  \label{fig:tournament}
\end{figure}

% ─────────────────────────────────────────────────────────────────────────────
\subsection{Guarded Tournament: Ranking Repairs Under Safety Gates}
\label{sec:tspo}
% ─────────────────────────────────────────────────────────────────────────────

Naively sampling inpainting hyperparameters wastes compute: many
configurations create artifacts, near-duplicates, or safety gains obtained by
destroying prompt fidelity. \textsc{GuardPaint} instead learns a tournament
policy that proposes repairs with high expected guarded utility.

\paragraph{Policy and state.}
The policy $\pi_\theta$ is a 3-layer MLP
($256\!\to\!128\!\to\!64$, SiLU, LayerNorm). It maps a compact repair state
to distributions over continuous knobs---CFG scale, mask dilation, mask
feather, noise jitter, inversion depth---and a discrete seed bucket:
\begin{align*}
s
=
\mathrm{concat}\Big[
&\mathrm{proj}(e_{\mathrm{text}}),\;
 \mathrm{proj}(z_t),\;
 \mathrm{proj}(e_{\mathrm{img}}),                                  \\
&\mathrm{proj}(\bar m),\;
 t/T
\Big]
\in\mathbb{R}^{257}.
\end{align*}
Each projection is 64-dimensional, $e_{\mathrm{text}}$ comes from the frozen
auditor BiLSTM, $e_{\mathrm{img}}$ is the auditor image embedding, and
$\bar m$ is the mask coverage fraction.

\paragraph{Candidate scoring.}
For a flagged region $(R,m)$, the policy samples $N=5$ configurations
$\{a_i\}_{i=1}^{N}$. The inpainter produces repairs $\{C_i\}_{i=1}^{N}$,
which are composed into the audit view and scored by $\As$:
\begin{align*}
(S_i,F_i,P_i,B_i)
&=
\As\!\left(
p,\;
\mathrm{Compose}(I_t,C_i,R)
\right).
\end{align*}
Here $S_i$ is policy safety, $F_i$ is prompt faithfulness, $P_i$ is
seam/perceptual quality, and $B_i$ is adversarial suppression confidence.

\paragraph{Guarded utility.}
Each candidate is ranked by a safety gain gated by quality and fidelity:
\begin{align*}
u_i
&=
\underbrace{(S_i-S_0-\delta)^+}_{\text{safety gain}}
\cdot
\underbrace{\mathbf{1}[P_i\ge\tau_P]}_{\text{quality gate}}
\\[-0.1em]
&\quad\cdot
\underbrace{\mathbf{1}[F_i\ge\tau_F]}_{\text{fidelity gate}}
\cdot
\underbrace{B_i}_{\text{suppression confidence}} .
\end{align*}
$S_0$ is the unedited control score, $\delta=0.01$ is the improvement margin,
and $\tau_P$, $\tau_F$ are calibrated gates. The winning edit is
$i^*=\arg\max_i u_i$ and is accepted if $u_{i^*}>0$. Otherwise, the
unedited control is retained. Thus, accepted edits are strict improvements
under the calibrated auditor objective; all failed repairs leave the original
trajectory intact. We use \emph{time-varying} thresholds that change
as a function of the normalized timestep $t_\text{norm} = t / T$:

\begin{equation}
\tau_P(t_{\text{norm}}) = 0.40 + 0.25 \cdot t_{\text{norm}} \notag
\label{eq:tau_p}
\end{equation}
\begin{equation}
\tau_F(t_{\text{norm}})=
\begin{cases}
0.30 + 0.30\, t_{\text{norm}}, \\
\quad t_{\text{norm}} < 0.85, \\[8pt]
0.55 - 0.10\,
\left(
\frac{t_{\text{norm}} - 0.85}{0.15}
\right), \\
\quad t_{\text{norm}} \geq 0.85.\notag
\end{cases}
\label{eq:tau_f2}
\end{equation}

\paragraph{Tournament policy objective.}
Utilities are converted into leave-one-out softmax credits:
\begin{align*}
w_i
&=
\mathrm{softmax}
\left(
\frac{u_i}{\tau}
\right)
-
\frac{1}{N},
\end{align*}
\begin{align*}
\tau=
\mathrm{std}\!\left(\{u_i\}_{i=1}^{N}\right).
\end{align*}
The offline policy objective is
\begin{align*}
\mathcal{L}_{\mathrm{tour}}
&=
-\sum_{i=1}^{N}
w_i\log\pi_\theta(a_i\mid s)
-
\lambda_H H[\pi_\theta(\cdot\mid s)]                                \\[-0.1em]
&\quad+
\lambda_c\,\mathrm{Cost}(\{a_i\}_{i=1}^{N})
-
\lambda_{\mathrm{div}}
\sum_{i<j}d(C_i,C_j).
\end{align*}
The entropy term prevents mode collapse, the cost term discourages expensive
settings such as deep inversion or excessive CFG, and the diversity term
encourages distinct candidates. We use
$\lambda_H^{\mathrm{cont}}=0.01$,
$\lambda_H^{\mathrm{disc}}=0.005$,
$\lambda_c=0.005$, and $\lambda_{\mathrm{div}}=0.01$, with mini-batches of
4 rollouts and gradient clipping at 1.0.
\iffalse
\paragraph{Listwise judge refinement.}
To maintain calibration between the tournament policy and the auditor’s
ranking signal, we co-train a lightweight
\texttt{ListwiseJudgeRefiner} that updates only the auditor’s
cross-attention and projection layers while keeping the ResNet-101
backbone frozen. Given tournament utilities $\{u_i\}_{i=1}^{N}$ and
auditor scores $\{s_i\}_{i=1}^{N}$, the target and model ranking
distributions are defined as
\begin{align*}
q_i = \frac{\exp(u_i/\tau_{\mathrm{target}})}{\sum_j \exp(u_j/\tau_{\mathrm{target}})},
\end{align*}
\begin{align*}
p_i = \frac{\exp(s_i/\tau_{\mathrm{model}})}{\sum_j \exp(s_j/\tau_{\mathrm{model}})},
\end{align*}
where $u_i$ is the utility of candidate $i$, $s_i$ is the auditor
safe score, $\tau_{\mathrm{target}}=0.5$ controls the
sharpness of the utility distribution, and
$\tau_{\mathrm{model}}=1.0$ controls the softness of the model
distribution. The refiner minimizes a combination of ListNet and top-1
ranking losses:
\begin{align*}
\mathcal{L}_{\mathrm{LN}} = -\sum_{i=1}^{N} q_i \log p_i,
\end{align*}
\begin{align*}
\mathcal{L}_{\mathrm{top1}} = \mathrm{CrossEntropy}(\mathbf{s}, [i^*]),
\end{align*}
where $i^*=\arg\max_i u_i$ denotes the highest-utility candidate. The
final objective is
\begin{align*}
\mathcal{L}_{\mathrm{judge}} =
\lambda_{\mathrm{LN}}\mathcal{L}_{\mathrm{LN}}
+
\lambda_{\mathrm{PL}}\mathcal{L}_{\mathrm{top1}},
\end{align*}
with $\lambda_{\mathrm{LN}}=1.0$ and
$\lambda_{\mathrm{PL}}=0.5$.
\fi
% ─────────────────────────────────────────────────────────────────────────────
\subsection{Latent Reinsertion Across Diffusion Families}
\label{sec:reinsertion}
% ─────────────────────────────────────────────────────────────────────────────

After a repair is selected, it must be reintroduced into the active generation trajectory without disrupting the target latent distribution. To achieve this across distinct architectural families, the reinsertion mapping is strictly governed by the underlying base model's specific forward diffusion or flow-matching equations.

\paragraph{Null-text inversion for UNet-based models.}
For SD~1.5 and SDXL, we use null-text inversion~\cite{mokady2023nulltext}. A
per-step null-text embedding is optimized for 10 gradient steps
(lr $=0.01$) to reconstruct the winning inpaint from the live latent. The edit
is then blended only inside the feathered mask:
\begin{align*}
z_{t-1}
&=
(1-\alpha_m)\,z^{\mathrm{ctrl}}_{t-1}
+
\alpha_m\,z^{\mathrm{edit}}_{t-1}.
\end{align*}
Here $\alpha_m$ is the mask upsampled to latent resolution. This keeps the
repair on the base model's denoising manifold and avoids drift from naive VAE
re-encoding.

\paragraph{Flow-ODE reinsertion for flow-matching models.}
For SD~3.5 and FLUX.1, Null-text inversion schedules are inapplicable. We instead use the
rectified-flow relation~\cite{liu2022flow}:
\begin{align*}
z_t
&=
(1-t)z_0+t\epsilon,                                                  \\[-0.1em]
\hat z_t
&=
(1-t_{\mathrm{norm}})\hat z_0
+
t_{\mathrm{norm}}\epsilon .
\end{align*}
The selected repair is encoded to $\hat z_0$ and projected forward to the
current normalized time. At late audit steps
($t_{\mathrm{norm}}<0.25$), where the trajectory is nearly deterministic, we
directly blend $\hat z_0$ into the live latent. DDPM-noise blending and DDIM
inversion are implemented as additional reinsertion baselines and compared in
Appendix~\ref{app:reinsertion_strategies}.

\begin{figure*}[!ht]
  \centering
  \includegraphics[width=\linewidth]{ASR.png}
  \caption{
    \textbf{Change-from-baseline heatmap} under the JailBreakDiffBench
    protocol~\cite{jailbreakdiffbench2024}.
    Each cell shows the absolute change in , ASR
    ($\downarrow$), of \textsc{GuardPaint} relative
    to the undefended base model. Blue cells indicate improvement; red cells
    indicate regression. Rows are attack families; columns are model architectures. Baselines are sourced from JailBreakDiffBench.
  }
  \label{fig:heatmap}
\end{figure*}

%% ════════════════════════════════════════════════════════════════════════════
\section{Full Pipeline Algorithm}
\label{sec:algorithm}
%% ════════════════════════════════════════════════════════════════════════════

Algorithm~\ref{alg:tipai} summarizes one audited decoding step. The base model
first advances normally, producing a control latent
$z^{\mathrm{ctrl}}_{t-1}$ and audit view $I_t$. If the current timestep is not
audited, the control latent is returned. Otherwise, the auditor runs once. A
benign image continues at the cost of a single auditor pass. A flagged image
triggers mask construction, candidate proposal, guarded scoring, and latent
reinsertion only if the winning candidate has positive utility.

\paragraph{Audit onset.}
Coherent visual structure typically appears after roughly 70--80\% of the
denoising trajectory. Earlier audits are inefficient because decoded views are
noise-dominated; later audits focus compute where unsafe semantic structure is
visible and repairable.

\begin{algorithm}[h]
\small
\caption{\textsc{GuardPaint} Decoding Step at Timestep $t$}
\label{alg:tipai}
\begin{algorithmic}[1]
\Require Base model $\Phi_{\mathrm{base}}$, auditor $\As$, inpainter $\Ag$,
         frozen policy $\pi_\theta$, prompt $p$, latent $z_t$,
         thresholds $(\delta,\tau_P,\tau_F)$
\State $z_{t-1}^{\mathrm{ctrl}}, I_t \gets
       \Phi_{\mathrm{base}}(z_t,p,t),\;
       \mathrm{Dec}(z_{t-1}^{\mathrm{ctrl}})$
\If{$t \notin \mathcal{T}_{\mathrm{audit}}$}
  \State \Return $z_{t-1}^{\mathrm{ctrl}}$
\EndIf
\State $r_0 \gets \As(p,I_t)$ \Comment{one auditor pass}
\If{$r_0.\mathrm{adv\_prob}<0.40$ \textbf{and}
    $r_0.\mathrm{harm\_class}\notin\{\mathrm{nudity},\mathrm{violence}\}$}
  \State \Return $z_{t-1}^{\mathrm{ctrl}}$ \Comment{benign path}
\EndIf
\State $S_0 \gets r_0.\mathrm{policy\_safe}$
\State $m \gets \mathrm{BuildMask}(r_0.\mathrm{heatmap})$
\State $s \gets \mathrm{StateEncode}(p,z_t,r_0.\mathrm{img\_embed},\bar m,t/T)$
\State $\{a_i\}_{i=1}^{N} \gets \pi_\theta(s)$
\State $\{C_i\}_{i=1}^{N} \gets \Ag(I_t,m,p;\{a_i\})$
\For{$i=1,\ldots,N$}
  \State $(S_i,F_i,P_i,B_i) \gets \As(p,C_i)$
  \State $u_i \gets (S_i-S_0-\delta)^+
      \mathbf{1}[P_i\ge\tau_P]\mathbf{1}[F_i\ge\tau_F]B_i$
\EndFor
\State $i^* \gets \arg\max_i u_i$
\If{$u_{i^*}>0$}
  \State $z_{t-1} \gets
  \mathrm{Reinsert}(C_{i^*},z_{t-1}^{\mathrm{ctrl}},m,t/T)$
\Else
  \State $z_{t-1} \gets z_{t-1}^{\mathrm{ctrl}}$
  \Comment{auditor-defined non-regression}
\EndIf
\State \Return $z_{t-1}$
\end{algorithmic}
\end{algorithm}

\paragraph{Complexity.}
Each audited step begins with one auditor pass $C_{\As}$. For benign prompts,
this is the only extra cost:
\begin{align*}
\mathrm{Cost}_t
&=
C_{\As}
\qquad
\text{if no trigger fires.}
\end{align*}
When repair is triggered, the cost scales with mined regions and candidate
repairs:
\begin{align}
\mathrm{Cost}_t
&=
C_{\As}
+
K_t \Bigl(
    N\,C_{\Ag}
    + N\,C_{\As} \notag \\
&\qquad\qquad
    + C_{\mathrm{reinsert}} 
\Bigr), \notag
\end{align}
where $K_t$ is the number of mined regions, $N$ is the tournament size, and
$C_{\Ag}$ is one inpainting pass. Initial text and image embeddings are cached
and reused across candidates. The tournament policy reduces wasted compute by
biasing proposals toward high-utility configurations, lower inversion depth,
and fewer exhausted candidate sets.

\begin{figure*}[!ht]
  \centering
  \includegraphics[width=\linewidth]{abalation.png}
  \caption{
    \textbf{Ablation: number of audited steps ($X$) vs.\ candidates per tournament
    ($N$) across five architectures.}
    Mean ASR ($\downarrow$) is shown as a function of audited denoising steps
    for $N\in\{1,3,5,10\}$. ASR decreases with both $X$ and $N$, with
    diminishing returns beyond $N=5$ .Flow-matching models such as FLUX.1 converge faster because their trajectories are more
    deterministic.
  }
  \label{fig:ablation}
\end{figure*}
















%% ════════════════════════════════════════════════════════════════════════════
\section{Experimental Setup}
\label{sec:experiments}
%% ════════════════════════════════════════════════════════════════════════════

\subsection{Adversarial Attack Benchmarks}

We evaluate under the \textbf{JailBreakDiffBench}
protocol~\cite{jailbreakdiffbench2024} on five black-box prompt-space attack families, where the adversary has only query access to the model with no knowledge of parameters, gradients, or latent representations:
\paragraph{ \textbf{SneakPrompt}}~\cite{yang2023sneakprompt}: formulates jailbreak as an RL-guided search that jointly optimizes adversarial prompts for semantic similarity to harmful targets and successful safety-filter evasion.
\paragraph{
\textbf{MMA}}~\cite{mma2024}: leverages offline CLIP-based guidance to iteratively perturb or replace tokens while preserving original prompt semantics in embedding space.
  \paragraph{\textbf{PGJ}}~\cite{pgj2024}: uses ChatGPT-generated antonym concepts to guide embedding-space optimization of stealthy adversarial prompts that bypass safety filters.
  \paragraph{\textbf{DACA}}~\cite{daca2024}: employs LLM-driven semantic rewriting to substitute trigger words with contextually benign alternatives, producing prompts that are textually innocent but semantically adversarial.
  \paragraph{\textbf{RABell}}~\cite{rabell2024}: applies surrogate CLIP-based guidance to iteratively replace tokens while maintaining surface-level naturalness and filter evasion.

\subsection{Base Models}

We attach GuardPaint to five T2I architectures spanning UNet-based
and flow-matching designs:
   Stable Diffusion 1.5 (SD~1.5; \textit{runwayml}) \cite{rombach2022ldm}
  SDXL Base 0.9 (SDXL; \textit{stabilityai}) \cite{podell2023sdxl}
  Stable Diffusion 3.5 Medium (SD~3.5-M; \textit{stabilityai})\cite{esser2024sd3}
  Stable Diffusion 3.5 Large-Turbo (SD~3.5-LT; \textit{stabilityai})
  FLUX.1-dev (\textit{Black Forest Labs}; flow-matching) \cite{flux1dev2026}
The Policy finetuned inpainter ($\Ag$) always uses SD~1.5 Inpainting, regardless of the base model family.

%% ════════════════════════════════════════════════════════════════════════════
\section{Results}
\label{sec:results}
%% ════════════════════════════════════════════════════════════════════════════

\subsection{Main Results}

We evaluate under the \textbf{JailBreakDiffBench} protocol~\cite{jailbreakdiffbench2024},
which provides standardized adversarial prompt generation, baseline results, and evaluation
scripts for fair cross-system comparison; all baseline ASR values are taken directly from
that benchmark.

Table~\ref{fig:heatmap} reports absolute attack success rates for all five attack families
across four base architectures, before and after \guardpaint{} defense. Across all
model--attack pairs, \guardpaint{} reduces ASR without exception.
The gains are largest on UNet-based models: on SD~1.5, DACA drops from
8.25\% to 1.94\% and SneakPrompt from 7.00\% to 1.79\%; on SDXL, MMA is
reduced to near zero (5.25\%\,$\to$\,3.92\%). Flow-matching models
(SD~3.5M, SD~3.5-LT, FLUX.1-dev) begin from lower baselines owing to stronger
safety pretraining, yet \guardpaint{} still achieves consistent reductions---for
instance, SneakPrompt on FLUX.1-dev falls from 8.25\% to 6.99\% and
MMA from 4.25\% to 2.44\%.

Prompt fidelity and image quality remain stable across all defended configurations. AlignS and BLIP scores stay within $\pm$0.05 of the undefended baseline, confirming that inpainting repairs are confined to the flagged region and leave the semantically valid portions of the prompt intact.

Flow-matching models present a distinct failure mode: when an attack succeeds despite defense, auditing and inpainting at early denoising steps causes the base model to partially override the reinserted latent, producing blurred or incoherent regions.

% ─────────────────────────────────────────────────────────────────────────────
% FIGURE PLACEHOLDER — Additional Qualitative / Comparison
% ─────────────────────────────────────────────────────────────────────────────
\begin{figure}[t]
  \centering
\includegraphics[width=\linewidth]{base_vs_comp.png}
  \caption{
    \textbf{Qualitative comparison.}
    Each row shows an adversarial prompt alongside the undefended
    baseline output, and the GuardPaint output.The policy aware inpainter produces semantically coherent safe alternatives rather than
    blank refusals, preserving scene context outside the flagged region.
  }
  \label{fig:qualitative}
\end{figure}

%% ════════════════════════════════════════════════════════════════════════════
%% ════════════════════════════════════════════════════════════════════════════
\section{Ablation Study}
\label{sec:analysis}
We ablate the two primary search parameters: steps audited~($X$) 
and candidates per tournament~($N$), which directly govern the 
security--compute trade-off.
Figure~\ref{fig:ablation} reports ASR across an attack  as a function of $X \in \{2,4,6,8\}$ for 
$N \in \{1,3,5,7\}$, evaluated on all five base architectures.

ASR decreases monotonically with both $X$ and $N$ across all 
model families.
Diminishing returns are pronounced: the gain from 
$N\!=\!1 \to N\!=\!3$ substantially exceeds that from 
$N\!=\!5 \to N\!=\!7$, and similarly for $X$ beyond~ steps.
Flow-matching models (FLUX.1) converge faster in $X$ due to 
their more linear and deterministic trajectories, where a 
single late-stage audit step captures most of the safety benefit.
These results motivate our production setting of $N\!=\!5$ and 
auditing the final two denoising steps, which sits on the knee 
of the quality--latency Pareto curve across all tested 
architectures (marked on each panel in Figure~\ref{fig:ablation}).

%%════════════════════════════════════════════════════════════════════════════
\section{Conclusion}
\label{sec:conclusion}
%% ════════════════════════════════════════════════════════════════════════════
 
We presented GuardPaint , a decoding-time safety framework for T2I diffusion models inspired by speculative decoding in LLMs.
By monitoring the diffusion trajectory with a multi-task auditor, rewriting suspect regions with a BCO-aligned inpainter, and efficiently proposing winning configurations with the Tournamnet policy, The GuardPaint framework provides compliant alternatives to adversarial outputs rather than blank refusals---without modifying any base model weights.The framework is plug-and-play across seven T2I architectures spanning four generations of design: DDPM-UNet (SD~1.5, SDXL), flow-matching-UNet (SD~3.5), and flow-matching-transformer (FLUX.1).


\clearpage
\newpage

%% ════════════════════════════════════════════════════════════════════════════
%% ════════════════════════════════════════════════════════════════════════════
\section{Limitations}
\label{sec:limitations}
%% ════════════════════════════════════════════════════════════════════════════

\paragraph{Latency.}
On a single NVIDIA A6000 (48 GB), typical overhead is
$\sim$7--8 seconds per audited timestep with $N=5$ for
UNet-based architectures (SD~1.5, SDXL), where each of the
5 SD~1.5 inpainting passes takes approximately 1.4 seconds;
on FLUX.1-dev, flow-ODE reinsertion raises this to
$\sim$20 seconds per audited timestep. Averaged across all
audit steps, the expected per-step overhead is $\sim$2 seconds for UNet
architectures.

\paragraph{Single-family inpainter.}
$\Ag$ is always an SD~1.5 Inpainting checkpoint, regardless of the base
model family. While reinsertion bridges the resolution and latent-space
mismatch, a native inpainter for each base family would reduce the fidelity
cost and better preserve the stylistic properties of higher-capacity
generators such as FLUX.1.

\paragraph{Coarse policy taxonomy.}
The auditor and inpainter are trained on three crude categories:
$\{\mathrm{safe},\,\mathrm{nudity},\,\mathrm{violence}\}$.
This taxonomy is insufficient for subtler policy violations that do not
manifest as localized explicit content. Harmful stereotyping, racial or
gender bias encoded in visual representations, cultural insensitivity,
and discriminatory imagery are not captured by any of the three training
labels and would not trigger the auditor regardless of severity.
Extending \textsc{GuardPaint} to such harms would require richer label
ontologies, dedicated training data, and potentially separate auditor
heads---each with their own calibration challenges.

\paragraph{Adversarial attacks on the auditor itself.}
\textsc{GuardPaint} assumes the auditor is a reliable detector within its
training distribution. A white-box or adaptive adversary with knowledge of
the auditor's architecture could craft inputs that simultaneously satisfy the
base model's generation objective and suppress the auditor's adversarial
score, bypassing the repair trigger entirely. Because the auditor operates in
image space and the inpainter operates in latent space, the two modalities
provide some implicit robustness---an adversary must fool both to achieve a
clean bypass---but this does not constitute a certified defense. Adaptive
attacks specifically targeting the auditor trigger threshold or the tournament
utility function remain an open threat.

\paragraph{Inpainting domain gap.}
The policy-aligned inpainter is fine-tuned on masks derived from auditor
heatmaps over a fixed training corpus. At inference, heatmap quality degrades
on out-of-distribution content styles---highly stylized art, photorealistic
renders with unusual lighting, or dense crowd scenes---which can produce
masks that are either too coarse (capturing safe context) or too sparse
(missing the unsafe region). In such cases the inpainter may repair the wrong
region, introduce visible seam artifacts, or fail the quality gate and leave
the base trajectory intact. The non-regression guard prevents active
degradation in the last scenario, but provides no safety improvement either.

\paragraph{Policy alignment without architectural diversity.}
Although $\Ag$ is preference-aligned via BCO to produce policy-compliant
repairs, alignment is constrained by the representational capacity of the
SD~1.5 inpainting backbone. Higher-capacity inpainters built on SDXL or
flow-matching architectures have richer generative priors that could
produce repairs with better perceptual quality, finer detail preservation,
and stronger semantic coherence with the surrounding unmasked context.
The resolution and latent-space mismatch that reinsertion currently
compensates for would not arise if the inpainter shared the same
architecture as the base generator.

\paragraph{Static calibration.}
Auditor thresholds $(\delta, \tau_P, \tau_F)$ and the tournament policy
$\pi_\theta$ are calibrated offline on a fixed distribution of prompts and
attack families. As jailbreak methods evolve, the proposal distribution
shifts and calibrated gates may become either too permissive or too
conservative. Periodic recalibration against new attack families is
necessary to maintain the intended operating point.

\paragraph{Semantic drift across reinsertion.}
Latent reinsertion---whether via null-text inversion for UNet models or
flow-ODE transport for flow-matching models---is an approximation. At early
audit steps where the trajectory is still noisy, reinsertion can introduce
subtle distributional drift that accumulates over subsequent denoising steps,
causing the final image to deviate from the intended prompt in ways that
neither the auditor nor the fidelity gate reliably detect. This effect is
most pronounced on long-horizon generation tasks with fine compositional
requirements.


\clearpage
\newpage

\bibliography{ref}


\clearpage
\newpage

\appendix
% ─────────────────────────────────────────────────────────────────────────────
\section{Training Dataset Construction}
\label{app:dataset}
% ─────────────────────────────────────────────────────────────────────────────

This appendix details the construction of the 82K image--prompt corpus used
to train the Auditor-Scorer $\As$, covering source datasets, prompt
reconstruction, labeling pipeline design, and labeling quality.

% ─────────────────────────────────────────────────────────────────────────────
\subsection{Source Datasets}
\label{app:dataset_sources}
% ─────────────────────────────────────────────────────────────────────────────

We draw from eight publicly available HuggingFace datasets spanning three
content categories: safe, nudity, and violence.
Table~\ref{tab:dataset} summarizes each source, its nominal category,
and whether prompts were available or reconstructed.

\begin{table*}[!ht]
\centering
\small
\setlength{\tabcolsep}{5pt}
\renewcommand{\arraystretch}{1.15}
\caption{Source datasets for the 82K auditor training corpus from huggingface.}
\label{tab:dataset}
\begin{tabular}{lcc}
\toprule
\textbf{Dataset} & \textbf{Category} & \textbf{Prompts} \\
\midrule
\texttt{Subh775/WeaponDetection} \cite{yolov7test_weapon_detection}              & Violence         & VLM-generated \\
\texttt{NeuralShell/Gore-Blood-Dataset-v1.0} \cite{Gore-Blood-Dataset-v1.0}  & Violence         & VLM-generated \\
\texttt{x1101/nsfw-full}                      & Nudity           & VLM-generated \\
\texttt{Lenkashell/unsafe\_violence\_image\_captions}  & Violence  & Existing \\
\texttt{Lenkashell/unsafe\_shocking\_image\_captions}  & Violence  & Existing \\
\texttt{yiting/UnsafeBench}\cite{10.1145/3719027.3765088}                   & Nudity, Violence & Existing \\
\texttt{OpenSafetyLab/t2i\_safety\_dataset}\cite{Li_2025_CVPR}   & Nudity, Violence & Existing \\
\bottomrule
\end{tabular}
\end{table*}

Sources were selected to cover both the nudity and violence policy axes with
sufficient volume and diversity of visual style, scene complexity, and
severity level. Datasets with pre-existing captions were used as-is after
cleaning; datasets without prompts required reconstruction (see
Section~\ref{app:prompt_reconstruction}).

% ─────────────────────────────────────────────────────────────────────────────
\subsection{Prompt Reconstruction}
\label{app:prompt_reconstruction}
% ─────────────────────────────────────────────────────────────────────────────

Four source datasets did not include generation prompts. For these, image
captions were synthesized using Qwen2.5-VL~\cite{bai2025qwen25vltechnicalreport}
as a surrogate prompt. The model was instructed to produce a concise,
descriptive caption in the style of a text-to-image generation prompt rather
than a natural-language description, targeting the visual content of the image
without editorializing about policy compliance. These reconstructed captions
serve the faithfulness branch of the auditor during training and are not used
for any safety label decision.

% ─────────────────────────────────────────────────────────────────────────────
\subsection{Labeling Pipeline}
\label{app:labeler_ablation}
% ─────────────────────────────────────────────────────────────────────────────

All images---regardless of source label---were relabeled into three mutually
exclusive categories: $\{\mathrm{safe},\,\mathrm{nudity},\,\mathrm{violence}\}$.
Multi-label annotation was evaluated and discarded: overlapping soft gradients
from multi-label supervision degraded auditor classification performance in
preliminary experiments. One-hot categorization was therefore adopted as the
final labeling scheme.

\paragraph{Labeler selection.}
We evaluated five candidate VLMs as automated labelers. Each was assessed on
a manually inspected reference set of hard cases including bikinis, underwear,
sports scenes with physical contact, crowds, close interpersonal proximity,
and classical artwork (e.g., Greek nude sculpture). The reference set was
chosen to stress-test policy-boundary behavior rather than gross violations,
which all models handle trivially.

Table~\ref{tab:labeler_ablation} summarizes observed failure modes.

\begin{table}[h]
\centering
\small
\setlength{\tabcolsep}{4pt}
\renewcommand{\arraystretch}{1.2}
\caption{Labeler ablation: observed failure modes on the hard-case calibration
set. FP = false positive rate; FN = false negative rate on the calibration
reference.}
\label{tab:labeler_ablation}
\begin{tabular}{p{0.22\linewidth}p{0.68\linewidth}}
\toprule
\textbf{Model} & \textbf{Observed failure mode} \\
\midrule
Qwen2.5-VL (7B) &
  Basic prompting yields high FP rate; chain-of-thought prompting introduces
  label contradictions; self-consistency is slow and inconsistent;
  few-shot prompting degrades prompt-following according to community
  reports. \\
\addlinespace
Qwen3-VL &
  Similar failure pattern to Qwen2.5-VL; no improvement on boundary cases. \\
\addlinespace
ShieldGemma (vanilla) &
  Slow inference; inconsistent categorization on borderline cases without
  fine-tuning; high FN rate overall. \\
\addlinespace
Grounding DINO~\cite{liu2023gdino} &
  Highly threshold-sensitive (small shift produces either high FP or high
  FN); lacks semantic context for borderline cases (\eg, classical nude
  sculpture). \\
\addlinespace
\textbf{InternVL-3.5 (8B)} &
  \textbf{Best overall performance}: lowest FP rate; strong contextual
  nuance (\eg, correctly categorizes classical nude sculptures as safe);
  conservative on FN (acceptable for training corpus construction). \\
\bottomrule
\end{tabular}
\end{table}

InternVL-3.5 (8B) was selected as the final labeler based on this ablation.

\paragraph{Why chain-of-thought prompting was de-emphasized.}
Explicit chain-of-thought prompting increases verbosity without reliably
improving category fidelity. On ambiguous images, it causes models to
generate unsupported post-hoc rationales or self-contradictory
intermediate reasoning that conflicts with the final label.
Self-consistency (repeated passes with majority voting) was also rejected: it
increases inference cost substantially while failing to correct systematic
classification errors caused by stable policy-boundary confusion---the primary
failure mode on borderline cases.

\paragraph{Policy-first prompt design.}
The InternVL-3.5 prompt was designed to be policy-oriented rather than
explanation-oriented. Categories are defined in operational terms directly tied
to visual evidence:

\begin{itemize}[leftmargin=1.2em,itemsep=2pt,topsep=2pt]
  \item \textbf{Nudity:} Swimwear, underwear,
        medical imagery, and athletic clothing are explicitly safe.
  \item \textbf{Violence:} visible physical harm or graphic cues such as
        blood or weapons in active use. Proximity, sparring poses, and
        non-graphic contact are safe.
  \item \textbf{Safe:} all other content, including classical artwork with
        nudity, sports imagery without graphic harm, and intimate proximity
        without explicit content.
\end{itemize}

This framing prevents surface-level cue triggering (e.g., visible skin area,
figure proximity, athletic posture) which was the dominant failure pattern
across rejected labelers.

% ─────────────────────────────────────────────────────────────────────────────
\subsection{Threshold Calibration and Dead-Zone Exclusion}
\label{app:threshold_calibration}
% ─────────────────────────────────────────────────────────────────────────────

InternVL-3.5 produces a continuous policy score for each image. Because no
large ground-truth labeled set was available for data-driven threshold
selection, the operating threshold was calibrated empirically against the
same manually inspected hard-case reference set used for labeler ablation.
Model scores were compared against human judgment on this subset, and the
threshold was set to reflect the desired moderation policy rather than an
externally optimized operating point.

A \textbf{dead zone} was reserved around the threshold: images whose scores
fell within a margin of the decision boundary were excluded from the training
corpus entirely rather than assigned an arbitrary label. This exclusion
strategy ensures that the high-confidence tails of the score distribution---
where label quality is highest---dominate training, while uncertain borderline
samples do not introduce spurious supervision signal.

\paragraph{Residual label noise.}
The most likely residual failure mode from single-shot VLM scoring is not
catastrophic random mislabeling but systematic boundary bias: InternVL-3.5
may mildly over-score skin-heavy but policy-safe images (e.g., bikini
photographs), causing a fraction of such cases to drift above the nudity
threshold. This produces moderate boundary noise in the minority class rather
than label collapse. For auditor training, such boundary noise is manageable
provided the high-confidence tails remain clean---which the dead-zone
exclusion strategy is designed to ensure.

% ─────────────────────────────────────────────────────────────────────────────
\subsection{Dataset Statistics and Splits}
\label{app:dataset_stats}
% ─────────────────────────────────────────────────────────────────────────────

The final labeled corpus contains 82,000 image--prompt pairs distributed
across the three categories. After dead-zone exclusion of boundary cases, the
corpus was split into training, validation, and test sets using a
70\,/\,15\,/\,15 stratified partition, preserving class proportions across
splits. The test set was held out entirely for the classification evaluation
reported in Table~\ref{tab:auditor_perf}.
\section{Multi-Modal Adversarial Auditor Pipeline Architecture}
\label{appendix:auditor_pipeline}

\begin{figure*}[t]
    \centering
    \includegraphics[width=\linewidth]{auditor-inference.png}
    \caption{
        \textbf{Architecture diagram of the Auditor module.}
    }
    \label{fig:auditor_architecture}
\end{figure*}

This appendix provides a rigorous architectural and mathematical specification
of the Adversarial Image Auditor pipeline. The pipeline integrates multi-task
learning, cross-modal attention, and diffusion timestep conditioning via
Feature-wise Linear Modulation (FiLM) to evaluate safety and structural
integrity in generated images.

% ─────────────────────────────────────────────────────────────────────────────
\subsection{Complete Pipeline Flow and Spatial Dimensions}

The system takes three distinct inputs: an image
$X \in \mathbb{R}^{B \times 3 \times 224 \times 224}$, a tokenized prompt
sequence $T \in \mathbb{R}^{B \times L}$ (where $L$ is the token sequence
length), and a normalized diffusion timestep $t \in \mathbb{R}^{B \times 1}$.
The primary data flow is formalised below.

\begin{enumerate}[leftmargin=*, label=\textbf{\arabic*.}]

\item \textbf{Visual Feature Extraction.}
\begin{equation}
\begin{aligned}
    F_{\text{visual}} = \text{ResNet101}_{\text{backbone}}(X)
    \;\in\; \\ \mathbb{R}^{B \times 2048 \times 7 \times 7} \notag
\end{aligned}
\end{equation}
A global average-pooled vector is also computed: 
\[
\begin{aligned}
f_{\text{global}} = \operatorname{AdaptiveAvgPool2d}(F_{\text{visual}}) \\ \in \mathbb{R}^{B \times 2048} \notag
\end{aligned}
\]

\item \textbf{Textual Encoding.}
\begin{equation}
    f_{\text{text}},\; S_{\text{seq}},\; M_{\text{pad}}
    = \text{BiLSTM}_{\text{encoder}}(T) \notag
\end{equation}
where $f_{\text{text}} \in \mathbb{R}^{B \times 512}$ is the global textual
context, $S_{\text{seq}} \in \mathbb{R}^{B \times L \times 512}$ contains
per-token sequence embeddings (pre-LayerNorm), and
$M_{\text{pad}} \in \{0,1\}^{B \times L}$ is a boolean padding mask.

\item \textbf{Prompt-to-Image Cross-Attention.}
Visual feature maps are projected to the textual dimension, flattened, and
normalised to form query vectors:
\begin{equation}
\begin{aligned}
    Q = \operatorname{LayerNorm}\!\bigl(\operatorname{Conv2d}_{1\times1}(F_{\text{visual}})\bigr) \\
    \;\in\; \mathbb{R}^{B \times 49 \times 512} \notag
\end{aligned}
\end{equation}
\begin{equation}
    K = V = \operatorname{LayerNorm}(S_{\text{seq}})
    \;\in\; \mathbb{R}^{B \times L \times 512} \notag
\end{equation}
Multi-head cross-attention yields prompt-conditioned visual representations:
\begin{equation}
\begin{aligned}
    A_{\text{attended}}
    = \operatorname{MultiHeadAttention}(Q,\,K,\,V;\; \\ \text{mask}=M_{\text{pad}})
    \;\in\; \mathbb{R}^{B \times 49 \times 512} \notag
\end{aligned}
\end{equation}
Spatial averaging gives $f_{\text{attended}} \in \mathbb{R}^{B \times 512}$,
which is used exclusively by the CLIP-style alignment head (Section~B.3).

\end{enumerate}

% ─────────────────────────────────────────────────────────────────────────────
\subsection{Timestep-Aware FiLM Conditioning}

Adversarial noise patterns manifest differently depending on the denoising
timestep $t \in [0,1]$.  A multi-layer perceptron maps $t$ to a shared
temporal embedding:
\begin{equation}
    e_{t} = \operatorname{MLP}_{\text{timestep}}(t)
    \;\in\; \mathbb{R}^{B \times 512} \notag
\end{equation}

This embedding drives two \emph{independent} FiLM projections with separate
learned weights --- one for adversary modulation and one for seam modulation.
Each projection produces scale $\gamma$ and shift $\beta$ coefficients for
Feature-wise Linear Modulation.

\begin{enumerate}[leftmargin=*, label=\textbf{\arabic*.}]

\item \textbf{Relative Adversary Modulation.}
The \emph{relative adversary score} quantifies the continuous perturbation
strength of an image on a $[0,1]$ scale (see Section~B.3).  To condition its
prediction on the timestep, we modulate the global pooled features:
\begin{equation}
\begin{aligned}
     \gamma_{\text{adv}},\;\beta_{\text{adv}}
    = \operatorname{Split}\!\bigl(\operatorname{Linear}_{\text{film,adv}}(e_t)\bigr) \\
    \;\in\; \mathbb{R}^{B \times 2048} \notag
    \end{aligned}
\end{equation}    

\begin{equation}
    f_{\text{global}}^{\text{mod}}
    = (1 + \gamma_{\text{adv}}) \odot f_{\text{global}} + \beta_{\text{adv}} \notag
\end{equation}
Here $\operatorname{Linear}_{\text{film,adv}} : \mathbb{R}^{512} \to
\mathbb{R}^{4096}$, and $\operatorname{Split}$ divides the output equally along
the last dimension.

\item \textbf{Seam Quality Modulation.}
A separate projection conditions the intermediate convolutional features
$F_{\text{seam}} \in \mathbb{R}^{B \times 512 \times 7 \times 7}$ used for
structural artefact detection:
\begin{equation}
\begin{aligned}
    \gamma_{\text{seam}},\;\beta_{\text{seam}}
    = \operatorname{Split}\!\bigl(\operatorname{Linear}_{\text{film,seam}}(e_t)\bigr) \\
    \;\in\; \mathbb{R}^{B \times 512} \notag
\end{aligned}
\end{equation}
\begin{equation}
    F_{\text{seam}}^{\text{mod}}
    = \bigl(1 + \gamma_{\text{seam}}^{\phantom{A}}\bigr) \odot F_{\text{seam}}
      + \beta_{\text{seam}} \notag
\end{equation}
where $\gamma_{\text{seam}}$ and $\beta_{\text{seam}}$ are broadcast over the
spatial dimensions.  Note that $\operatorname{Linear}_{\text{film,seam}}$ maps
$\mathbb{R}^{512} \to \mathbb{R}^{1024}$ and has no shared weights with
$\operatorname{Linear}_{\text{film,adv}}$.

\end{enumerate}

% ─────────────────────────────────────────────────────────────────────────────
\subsection{Specialised Multi-Task Auditor Heads}

The auditor decouples distinct tasks into dedicated prediction heads to ensure
stable gradient flow and to protect alignment encodings from direct safety-label
contamination.

\begin{itemize}[leftmargin=*]

\item \textbf{Binary Adversarial Head.}
Computes $P(\text{Adversarial} \mid X)$ directly from the \emph{unmodulated}
spatial features $F_{\text{visual}}$ via a $1{\times}1$ convolution, adaptive
global average pooling, and Sigmoid activation.  This head intentionally
bypasses FiLM conditioning so that the binary decision is not conflated with
timestep dynamics.

\item \textbf{Safety Category Head.}
Projects $F_{\text{visual}}$ to $C=3$ channels and applies cost-sensitive
cross-entropy with class weights $w = [1.0,\,5.0,\,2.0]$.  The elevated weight
for class~1 (adversarial) reflects its underrepresentation in the training
distribution; the moderate weight for class~2 (borderline) penalises
ambiguous predictions more than the majority safe class.

\item \textbf{Relative Adversary Score Head.}
Regresses the continuous perturbation strength $s \in [0,1]$ from the
FiLM-modulated vector $f_{\text{global}}^{\text{mod}}$ using a three-layer MLP
with dropout.  Unlike the binary head, this head is conditioned on the timestep
embedding to capture how perturbation magnitude varies across denoising stages.

\item \textbf{Seam Quality Assessment Head.}
Evaluates inpainting seams and composite boundaries by applying a regression
head to $F_{\text{seam}}^{\text{mod}}$, yielding a localised artefact score.

\item \textbf{CLIP-style Faithfulness Alignment Head.}
Projects $f_{\text{attended}}$ and $f_{\text{text}}$ independently into a
shared latent space $\mathbb{R}^{256}$ using separate two-layer MLPs, then
optimises image--prompt alignment via symmetric InfoNCE loss with a learnable
log-temperature $\tau$.  Both the cross-attended visual vector and the raw
global text vector $f_{\text{text}}$ are fed to this head to preserve low-level
lexical grounding alongside spatially-attended semantics.

\end{itemize}

% ─────────────────────────────────────────────────────────────────────────────
\subsection{GradCAM Explainability Pipeline}

For visual explainability, target activations are extracted from the final
residual convolutional layer of the ResNet101 backbone.  The localised risk
heatmap $M_{\text{risk}} \in \mathbb{R}^{7 \times 7}$ for safety class $c$ is:
\begin{equation}
    M_{\text{risk}}
    = \operatorname{ReLU}\!\left(\sum_{k} \alpha_k^c \cdot F_{\text{visual},k}\right) \notag
\end{equation}
where the importance weights $\alpha_k^c$ are obtained by global average
pooling of the backpropagated gradients:
\begin{equation}
    \alpha_k^c
    = \operatorname{AdaptiveAvgPool2d}\!\left(
        \frac{\partial\, y^c}{\partial\, F_{\text{visual},k}}
      \right) \notag
\end{equation}
The resulting maps are bilinearly upsampled to $224{\times}224$ and
alpha-blended onto the source image using custom contour overlays to localise
flagged regions.  Because GradCAM operates directly on $F_{\text{visual}}$, it
provides explanations that are independent of FiLM conditioning, giving an
unbiased view of which image regions drive the safety classification.


%% ════════════════════════════════════════════════════════════════════════════
%% APPENDIX: Policy-Aligned Inpainter --- SFT and BCO Training Pipeline
%% ════════════════════════════════════════════════════════════════════════════

\section{Policy-Aligned Inpainter: SFT and BCO Training Details}

This appendix documents the complete training pipeline for the policy-aligned
inpainter $\mathcal{A}_g$, covering: the Refusal SFT stage
(\S\ref{app:sft}), the alignment objective selection and the migration from KTO to BCO (\S\ref{app:algo_choice}--\ref{app:kto_to_bco}), the full BCO loss formulation and motivations (\S\ref{app:bco_full}), the
failure archaeology from all training runs (\S\ref{app:failure_modes}), and
the final validated training configuration (\S\ref{app:final_config}).

%% ────────────────────────────────────────────────────────────────────────────
\subsection{Alignment Objective Selection: Why KTO and BCO}
\label{app:algo_choice}
%% ────────────────────────────────────────────────────────────────────────────

\paragraph{Dataset constraints.}
Modern RLHF methods impose different data-preparation burdens.
DPO~\cite{wallace2023diffusion_dpo} requires \emph{paired} (chosen, rejected)
completions for every prompt; for images this demands either human preference
labelling or an expensive oracle model.

KTO~\cite{li2024diffusion_kto} and BCO~\cite{jung2024bco} both operate on
\emph{unpaired} binary-labelled data $(x, y)\in\{0,1\}$, where $y=1$
denotes a safe (desirable) completion and $y=0$ denotes a policy-violating
(undesirable) one.
This format is produced directly by the auditor, requiring no
additional human annotation.

\paragraph{Risk profile motivation.}
For a content-safety inpainter, the error costs are highly asymmetric:
a false negative (generating adversarial material undetected) is far more
costly than a false positive (over-suppressing a benign region).
KTO was initially chosen specifically for its loss-averse utility function
and its asymmetric treatment of gains and losses, properties that appeared
well-suited to this risk profile.

%% ────────────────────────────────────────────────────────────────────────────
\subsection{Refusal SFT - LoRA Fine-Tuning}
\label{app:sft}
%% ────────────────────────────────────────────────────────────────────────────

\paragraph{Motivation: the know-before-you-refuse principle.}
The LLM training paradigm provides a useful analogy: instruction tuning
precedes alignment because a model must first \emph{know how to perform a
task} before it can refuse it in a principled way.
Applying an alignment signal directly to a model with no prior concept of
``safe inpainting'' is under-constrained: the model has no learned attractor
basin for policy-compliant completions, so the alignment gradient has
nothing to guide it toward.

Formally, let $\pi_\theta$ be the inpainter policy and $\mathcal{R}$ the
target ``refusal manifold'' in latent space---the set of completions that
replace NSFW content with spatially coherent, policy-compliant alternatives.
An alignment objective such as BCO applies a gradient
$\nabla_\theta \mathcal{L}_\text{BCO}$ that rewards distance from the
policy-violating attractor.
However, if $\mathcal{R}$ is not already in the support of $\pi_\theta$,
the gradient has no preferred direction in which to move the policy; it
merely pushes it \emph{away} from the violation without specifying
\emph{where} to go.
SFT pre-seeds the policy with a collapsed distribution over $\mathcal{R}$,
giving the BCO gradient a meaningful signal basin to expand.

In our training logs, this manifests clearly: the BCO run starting from an
SFT-initialised checkpoint exhausts its SFT prior through approximately
step~500, then searches the manifold for an exploit.
The asymmetric auxiliary losses (reconstruction and identity, detailed
in \S\ref{app:bco_full}) constrain this search to a corridor that
eventually leads to the globally optimal solution---rendering coherent
clothing.
A BCO run starting from a vanilla SD-inpainting checkpoint does not exhibit
this three-phase structure; it collapses directly to the pixel-smudging
attractor with no recovery.

\paragraph{Base model and adapter.}
SFT training starts from \emph{runwayml/stable-diffusion-inpainting}, a
9-channel inpainting UNet conditioned on
$(z_t \in \mathbb{R}^{4}, m_\ell \in \mathbb{R}^{1}, \tilde{z}_0 \in
\mathbb{R}^{4})$, where $m_\ell$ is the latent-space mask and $\tilde{z}_0$
is the masked image latent.
We apply LoRA~\cite{hu2022lora} with rank $r{=}64$, $\alpha{=}64$
(effective scale $= \alpha/r = 1.0$, i.e.\ no implicit scaling)
targeting the following parameter groups:
\begin{itemize}[leftmargin=1.5em,itemsep=1pt,topsep=2pt]
  \item Cross- and self-attention:
        \emph{to\_q}, \emph{to\_k}, \emph{to\_v}, \emph{to\_out.0}
  \item Feed-forward layers in transformer blocks:
        \emph{ff.net.0.proj}, \emph{ff.net.2}
  \item Spatial projection layers: \emph{proj\_in}, \emph{proj\_out}
  \item ResNet convolutional layers (structural/spatial output):
        \emph{conv1}, \emph{conv2}
  \item UNet input convolution layer: \texttt{conv\_in}
\end{itemize}

The inclusion of \emph{conv\_in} is critical.
The 9-channel inpainting UNet differs from the standard 4-channel text-to-image
UNet in precisely this first layer: it processes the concatenated
$(z_t, m_\ell, \tilde{z}_0)$ tensor.
Leaving \emph{conv\_in} frozen during alignment causes the
mask-conditioning logic to resist policy signals, because the gradient
from the loss cannot propagate into the layer that interprets which pixels
to modify.
Similarly, the ResNet convolutional layers (\emph{conv1}, \emph{conv2})
handle the structural and spatial content of the output.
Refusal requires rendering new content (clothing, blurred faces)---a
structural operation---and attention-only LoRA lacks the representational
capacity for such changes.

Stage B-2 uses rank $r{=}128$, $\alpha{=}128$ for the alignment phase,
on the hypothesis that the earlier runs with smaller rank were exhausting
LoRA capacity and saturating prematurely.

\paragraph{Training data.}
The SFT dataset consists of approximately 2\,000 images (1\,000 nudity
violations, 1\,000 violence violations).
For each unsafe source image, the auditor generates a mask over the
policy-violating region.
The source image and mask are then passed to the inpainter with a
custom-generated safe prompt (derived by prompting a Qwen-2.5-VL captioner
followed by a LLaMA-8B rewriter) to produce a target ``accepted'' image.
Training pairs are therefore $(x_\text{unsafe}, m, x_\text{safe}, p_\text{unsafe})$:
the unsafe source, its violation mask, the human-approved safe target, and
the original adversarial prompt.

Training on the \emph{adversarial} prompt rather than the safe prompt is
deliberate: at inference, the inpainter receives the adversarial prompt
because the user's intent is unknown.
The model must learn the pixel-space signature of a safe completion
conditioned on an unsafe prompt---not relying on text-based cues to trigger
refusal.

\paragraph{Loss and schedule.}
The SFT loss is Min-SNR~\cite{hang2023minsnr} weighted MSE across the full
noise schedule $t \in [0, 1000)$:
\[
  \mathcal{L}_\text{SFT} = \mathbb{E}_{t,\epsilon}\!\left[
    w(t)\cdot
    \|\epsilon_\theta(\begin{aligned}z_t, t, m_\ell, \\ \tilde{z}_0, p_\text{unsafe} \end{aligned}) 
     - \epsilon\|^2
  \right],
  \label{eq:sft_loss}
\]
where $w(t) = \min(\text{SNR}(t), \gamma)/\text{SNR}(t)$ with $\gamma{=}5$.
Without Min-SNR reweighting, high-noise timesteps (large $t$) dominate the
gradient because the signal-to-noise ratio is low and the loss magnitude is
large.
This matters here because safe completions are visually coherent only at
\emph{low} noise levels ($t \ll 500$); high-noise gradients contribute
primarily noise-level statistics and not the structural pattern of
clothing or background fill.
Min-SNR redistributes gradient mass toward the low-noise regime where the
visual content of the refusal is actually resolved.

Formally, the Min-SNR weight is derived from the signal-to-noise ratio
$\text{SNR}(t) = \bar\alpha_t / (1 - \bar\alpha_t)$, where
$\bar\alpha_t = \prod_{s=1}^t \alpha_s$ is the cumulative product of the
noise schedule.
The clamp at $\gamma$ prevents the weight from collapsing to zero at very
low $t$ (where SNR is enormous), maintaining gradient flow throughout the
schedule.

A small noise offset of $0.05$ is added:
$\epsilon' = \epsilon + 0.05 \cdot \epsilon_0$ where $\epsilon_0 \sim
\mathcal{N}(0, I_{C \times 1 \times 1})$ is broadcast.
This offset helps with dark and saturated regions---common in clothing
inpainting---by preventing the VAE latent distribution from collapsing into
degenerate attractors near the boundary of its quantization grid.

After convergence (loss plateau below $\sim$0.07 over 8 epochs), the LoRA
weights are merged into the base checkpoint via \texttt{merge\_and\_unload()},
producing a single 9-channel UNet used as the Stage B-2 initialisation.

\paragraph{How SFT initialises BCO: a geometric argument.}
Let $\mathcal{M}_\text{img}$ denote the natural image manifold in latent
space and $\mathcal{M}_\text{safe} \subset \mathcal{M}_\text{img}$ the
submanifold of safe completions inside the mask.
Pre-SFT, $\pi_\theta$ (the base inpainting model) has learned a strong
prior toward $\mathcal{M}_\text{img}$ conditioned on the input context, but
has never encountered a training signal that distinguishes
$\mathcal{M}_\text{safe}$ from the rest of the manifold.

SFT provides approximately 2\,000 examples of $(x_\text{unsafe}, m,
x_\text{safe})$ triples.
The MSE loss \eqref{eq:sft_loss} carves a narrow basin in the loss
landscape centred on $\mathcal{M}_\text{safe}$ for inputs of the form
$(z_t, m_\ell, \tilde{z}_0, p_\text{unsafe})$.
After SFT, $\pi_\theta$ has a collapsed mode at clothing/covering
completions for masked unsafe inputs.

BCO then operates on this initialised policy.
The BCO gradient moves the policy such that:
(a) on safe inputs, the policy stays near or improves upon the reference;
(b) on unsafe inputs, the policy deliberately degrades relative to the
reference (moves further from the policy-violating content).
Because the SFT basin already exists, the BCO search for ``worse on
unsafe'' lands directly on $\mathcal{M}_\text{safe}$ rather than on
degenerate off-manifold solutions.
This explains the three-phase dynamics observed in BCO training: the SFT prior
is exploited first (Phase 1), the degenerate attractor is blocked by
auxiliary losses (Phase 2), and the SFT basin is rediscovered as the
unique viable optimum (Phase 3).

%% ────────────────────────────────────────────────────────────────────────────
\subsection{KTO: Formulation, Experiments, and Failure Modes}
\label{app:kto_to_bco}
%% ────────────────────────────────────────────────────────────────────────────

We now document the KTO implementation and the concrete failure modes that motivated the switch to BCO.

\subsubsection{KTO Formulation}

The KTO objective~\cite{ethayarajh2024ktomodelalignmentprospect} for diffusion models \cite{li2024diffusion_kto} adapts the
Kahneman--Tversky prospect-theoretic value function to the noise-prediction setting.
Our implementation computes a ``g-term'' as:
\begin{equation}
  g_i = \ell^\theta_i - \ell^\text{ref}_i, \notag
  \label{eq:kto_gterm}
\end{equation}
where $\ell^\bullet_i$ is the masked MSE of model $\bullet$ against the
clean latent $z_0$ inside the mask:
\begin{equation}
  \ell^\bullet_i
  = \frac{\sum_{b,c,h,w} (1 + \lambda_m m_\ell) \cdot
      (\hat{z}_0^\bullet - z_0)^2 \cdot m_\ell}
    {\sum m_\ell}. \notag
  \label{eq:masked_mse}
\end{equation}
A positive $g_i$ means the trainable model has drifted further from $z_0$
than the reference inside the mask---\emph{worse} alignment with the safe target.

The signal is centered using an EMA over safe samples:
\[
  \delta_\text{ema} \leftarrow 0.99\,\delta_\text{ema}
    + 0.01\,\bar{g}|_{\text{safe}},
  \quad \label{eq:kto_ema}
\]
\[
  \delta_\text{ema} \in [-0.02, 0.05].
\]
The centered signal drives a sigmoid-based loss:
\begin{equation}
  h_i = \sigma\!\left(\text{sgn}_i \cdot \beta \cdot (g_i - \delta_\text{ema})\right),
  \quad \notag
  \end{equation}
  \begin{equation}
  \mathcal{L}_\text{KTO} = \mathbb{E}_i\!\left[w_y \cdot (1 - h_i)\right], \notag
  \label{eq:kto_loss}
\end{equation}
where $\text{sgn}_i = +1$ for safe samples and $-1$ for unsafe,
$\beta = 7$, and $w_y \in \{1.0, 5.0, 12.0\}$ for safe, nudity, and
violence classes respectively.

\subsubsection{KTO Experiment Progression}

\paragraph{Vanilla KTO.}
The first KTO implementation used a static reference model and the
$\epsilon$-space MSE as a proxy for the probability of generating
policy-violating content.
The model quickly found a degenerate attractor: generating uniform blue
patches inside the masked region.
This trivially reduced the noise-prediction MSE difference relative to the
reference (whose MSE was also high at the same patches) while producing
visually incoherent output.
The model had hacked the reward by leaving the image manifold entirely.

\paragraph{Dynamic reference.}
To prevent manifold departure, we tried updating the reference model every
$N$ gradient steps.
This accelerated instability: accumulated errors in the trainable model
were passed to the reference model, which then guided training toward its
own failures in a positive-feedback loop.
Image quality outside the mask degraded correspondingly.

\paragraph{Reconstruction and identity losses.}
To anchor the model to the image manifold, we introduced two auxiliary
losses (described fully in \S\ref{app:bco_full}):
a reconstruction loss penalising noise-prediction error outside the mask,
and an identity guardrail penalising $\hat{z}_0$-space divergence from the
reference.
These improved manifold stability but introduced a new failure mode: the
model learned to smudge hands, feet, and peripheral body parts inside the
mask---a pixel-destruction strategy that achieved low KTO/BCO loss while
staying near the identity boundary.

\paragraph{Modified loss formulation.}
We observed that the reference model did not always generate a
policy-violating inpaint, especially for prompts that did not explicitly
mention NSFW content (Qwen-2.5-VL captions tend to be descriptive rather
than explicit).
We reframed the problem as alignment \emph{and} safety-quality improvement,
computing the g-term as the difference of MSEs against ground-truth noise
rather than against the reference.
Results were similar.

\paragraph{Prompt dropout.}
Observing that the cross-attention layers were overweighted (higher
activation delta than before the manifold dip and after recovery), we
suspected the cross-attention was learning to key on NSFW token patterns
rather than visual features.
We introduced 30\% prompt dropout per batch, zeroing the text conditioning
to force spatial grounding of the refusal signal.
This engineering choice is retained in the BCO configuration.

\subsubsection{Why KTO Failed: Theoretical Analysis}

\paragraph{Failure Mode 1: $\epsilon$-space stochasticity.}
The KTO g-term \eqref{eq:kto_gterm} is computed from $\hat{z}_0$
predictions, but the $\hat{z}_0$ reconstruction is itself a function of
the noise $\epsilon$ sampled at each step.
At high timesteps $t \sim \mathcal{U}[500, 1000)$, the standard deviation
of $\epsilon$ is $\sim\!\mathcal{O}(1)$, and the resulting variance in the
masked MSE across samples from the same image at the same $t$ is
$\mathcal{O}(10^{-1})$.
The preference gap between safe and unsafe images in our dataset is
$\mathcal{O}(10^{-2})$.
Consequently, the signal-to-noise ratio of the preference gradient is less
than 1; the loss tracks noise-schedule variance rather than safety
information.
This produces the characteristic flat
$h_\text{safe} \approx h_\text{unsafe} \approx 0.500$ plateau observed in
early KTO runs.

\paragraph{Failure Mode 2: Asymmetric risk cannot be encoded cleanly.}
KTO uses $(1 - \sigma(x))$ as its loss function, which is equivalent to
$\log \sigma(-x)$ up to sign.
The curvature of this function is symmetric in $x$: moving the policy away
from a violation provides the same gradient magnitude as moving it toward
a safe completion.
Our risk profile is strongly asymmetric: false negatives on unsafe content
(NSFW output reaches the user) are far more costly than false positives
(over-suppression of benign content).
KTO has no mechanism to encode this asymmetry beyond the class weights
$w_y$, which scale the loss magnitude but not the gradient shape.

Additionally, KTO's EMA baseline is computed from safe samples only
(\eqref{eq:kto_ema}), then subtracted from both safe \emph{and} unsafe
g-terms.
This conflates the distributional statistics of two distinct populations:
safe images (where $g \approx 0$ under a well-behaved policy) and unsafe
images (where $g \gg 0$ is desirable).
Subtracting a safe-derived baseline from the unsafe signal systematically
miscalibrates the centering.

\paragraph{Failure Mode 4: Identity loss bleeds globally.}
The KTO implementation computes the identity gap as a single global scalar:
\begin{equation}
  \ell_\text{id} = \big(\|{\hat{z}_0^\theta - \hat{z}_0^\text{ref}}\|^2 - 0.02\big)_+^2, \notag
\end{equation}
aggregated over all samples.
A single severely drifted sample inflates this scalar, raising the penalty
for all other samples and causing the optimizer to over-regularize the
well-behaved ones.
The resulting signal is poorly calibrated and impedes the policy's ability
to diverge from the reference on unsafe samples---which is precisely the
desired behavior.

%% ────────────────────────────────────────────────────────────────────────────
\subsection{BCO Loss: Full Formulation with Proofs and Motivations}
\label{app:bco_full}
%% ────────────────────────────────────────────────────────────────────────────

BCO~\cite{jung2024bco} reformulates the alignment objective as binary
classification: a sample is ``desirable'' ($y=1$) if it should be
encouraged and ``undesirable'' ($y=0$) if it should be suppressed.
The following sections derive each component of the BCO loss from first
principles.

\subsubsection{Why $\hat{z}_0$, Not $\epsilon$}

Standard diffusion training minimises noise-prediction MSE
$\|\epsilon_\theta - \epsilon\|^2$.
At timestep $t$, the magnitude of $\epsilon$ is $\sim\!\mathcal{O}(1)$, so
small changes in $\epsilon_\theta$ produce large, timestep-dependent changes
in the MSE.
A reward computed from $\|\epsilon_\theta - \epsilon\|^2$ is therefore
dominated by $t$, making per-class weighting meaningless---a sample at
$t=900$ produces 10--100$\times$ the loss magnitude of the same sample at
$t=100$.

Instead, both UNets predict the clean latent:
\begin{equation}
  \hat{z}_0^\bullet
  = \frac{z_t - \sqrt{1-\bar\alpha_t}\,\epsilon_\bullet}{\sqrt{\bar\alpha_t}},
  \quad \bullet \in \{\theta, \text{ref}\}. \notag
  \label{eq:z0hat}
\end{equation}
The denominator $\sqrt{\bar\alpha_t}$ ranges from approximately $1.0$
(at $t=0$) to $\sim$0.01 (at $t=999$).
While $\hat{z}_0$ predictions are noisier at high $t$, they are calibrated
to the same latent space regardless of $t$: $\hat{z}_0 \in
[-3\sigma_{z_0}, +3\sigma_{z_0}]$ with bounded variance.
This invariance is the key property: per-class weights now operate on a
reward signal with consistent scale, making nudity weight $5.0\times$ and
violence weight $12.0\times$ interpretable as actual preference ratios.

\subsubsection{Masked MSE Reward}

The reward quantifies how much better the trainable policy reconstructs
the safe target inside the mask, relative to the frozen reference:
\begin{align}
  w &= 1 + \lambda_m \cdot m_\ell, \quad \lambda_m = 0.5, \notag
  \\
  \ell_\bullet &= \frac{\sum_{b,c,h,w} w \cdot (\hat{z}_0^\bullet - z_0)^2 \cdot m_\ell}
                       {\sum m_\ell}, \notag
                       \label{eq:masked_mse_bco} \\
  g_i &= \ell^\text{ref}_i - \ell^\theta_i. \notag 
\end{align}
Note the sign convention: $g_i > 0$ means the trainable model reconstructs
the masked region \emph{better} than the reference (closer to $z_0$).
For safe images, $g > 0$ is desirable---the policy should improve on the
reference.
For unsafe images, $g < 0$ is desirable---the policy should produce
reconstructions \emph{further} from the policy-violating $z_0$.

The mask weight $\lambda_m = 0.5$ applies $1.5\times$ gradient emphasis
inside the adversarial region without completely ignoring the surrounding
context.
Ignoring the unmasked region entirely ($\lambda_m \to \infty$, i.e.\ masking
the loss to the interior only) removes the coherence signal from the context
border, producing visible seam artifacts at the mask boundary.

\paragraph{Hinge cap on unsafe reward.}
For unsafe samples ($y=0$), the reward is capped at:
\begin{equation}
  g_i \leftarrow \min\!\Big(g_i,\;
    1.5 \cdot |\bar{g}_\text{unsafe}|_\text{detach}\Big), \notag
  \label{eq:hinge_cap}
\end{equation}
where $\bar{g}_\text{unsafe}$ is the mean reward over unsafe samples in
the current batch.
Without this cap, the policy can exploit the mask by uniformly suppressing
\emph{all} content inside it---a trivial maximisation of $-g$ that
achieves a very negative reward (far from $z_0$) while producing visually
incoherent output.
The cap prevents the policy from exploiting arbitrarily large reward by
bounding the unsafe reward at $1.5\times$ its own class mean,
forcing it to find solutions that are both safe-suppressing \emph{and}
spatially coherent.

\subsubsection{Reward-Shift EMA and BCO Loss}

To center the reward at the decision boundary between safe and unsafe
classes, we maintain an exponential moving average:
\begin{align}
  \delta_\text{raw} &= \tfrac{1}{2}
    \Big(\mathbb{E}[g | y{=}1] + \mathbb{E}[g | y{=}0]\Big), \notag
    \label{eq:delta_raw} 
\end{align}
\[
  \delta_\text{ema} \leftarrow
    0.999\,\delta_\text{ema} + 0.001\,\delta_\text{raw}, \label{eq:delta_ema_1}
\]
\[
  \quad \delta_\text{ema} \in [-0.03, +0.03]. \label{eq:delta_ema_2}
\]

\paragraph{Why the symmetric midpoint?}
The shifted reward $r_i = g_i - \delta_\text{ema}$ should be positive for
safe samples (the policy is better than baseline) and negative for unsafe
samples (the policy is worse than baseline).
The optimal centering point is the midpoint of the two class means, which is
exactly $\delta_\text{raw}$.
KTO's centering (\eqref{eq:kto_ema}) uses a safe-only EMA, which
systematically underestimates the true decision boundary when $g$ differs
across classes, miscalibrating the loss.

The EMA clamp is applied \emph{in-place} to the stored EMA value before any downstream use:
The $[-0.03, +0.03]$ bound is tighter and symmetric compared to KTO's $[-0.02, +0.05]$, preventing runaway baseline drift in either direction.

\paragraph{BCO loss: softplus vs.\ sigmoid.}
Given $r_i = g_i - \delta_\text{ema}$ and label sign
$\text{sgn}_i = 2y_i - 1 \in \{-1, +1\}$, the per-sample BCO loss is:
\begin{equation}
\begin{aligned}
  \ell^\text{BCO}_i = \mathrm{softplus}(-\,\text{sgn}_i \cdot \beta \cdot r_i)
  \\= \log(1 + \exp(-\,\text{sgn}_i \cdot \beta \cdot r_i)). \notag
  \label{eq:bco_loss}
\end{aligned}
\end{equation}

\paragraph{Why softplus instead of KTO's $(1 - \sigma)$ loss?}
Note that $\log \sigma(x) = -\text{softplus}(-x)$.
The KTO per-sample loss $w_y(1 - h_i)$ with $h_i = \sigma(\cdot)$ is
related to the negative log-likelihood $-\log \sigma(\cdot)$ but has a
different gradient profile.

\begin{enumerate}[leftmargin=1.5em,itemsep=2pt]
  \item \textbf{Gradient saturation.}
  $\partial(1-\sigma(x))/\partial x = -\sigma(x)(1-\sigma(x))$, which
  collapses to zero as $x \to \pm\infty$.
  $\partial\,\mathrm{softplus}(-x)/\partial x = -\sigma(-x) = -(1-\sigma(x))$,
  which goes to $-1$ as $x \to +\infty$ and to $0$ as $x \to -\infty$.
  Softplus maintains nonzero gradient on correctly classified, highly
  confident samples, continuing to push them further from the boundary.
  KTO's $(1-\sigma)$ loss provides a vanishing gradient once $h$ is large,
  causing training to stall for well-separated samples while allowing
  hard cases to drift.

  \item \textbf{Probabilistic interpretation.}
  Softplus is the numerically stable form of the binary cross-entropy
  loss $-\log \sigma(\cdot)$, which is the maximum-likelihood loss for a
  logistic classifier.
  KTO's $(1-\sigma)$ loss lacks this clean probabilistic interpretation
  and is less aligned with standard RL reward maximization intuition.

  \item \textbf{Numerical stability.}
  At $\beta=50$, the argument to the sigmoid can easily reach $|\beta r|
  \sim 50 \times 0.5 = 25$, causing floating-point overflow in
  $\exp(25) \approx 7 \times 10^{10}$.
  PyTorch's \texttt{F.softplus} handles this via the numerically stable
  formulation $\max(x, 0) + \log(1 + \exp(-|x|))$.
\end{enumerate}

The final BCO loss aggregates with per-class weights $w_y$:
\begin{equation}
\begin{aligned}
    \mathcal{L}_\text{BCO} = \mathbb{E}_i\!\left[
    w_y \cdot \mathrm{softplus}(-\,\text{sgn}_i \cdot \beta \cdot r_i)
  \right], \\
  \quad \beta = 50.
  \notag
  \label{eq:bco_loss_final} 
\end{aligned}
\end{equation}

\subsubsection{Identity Guardrail: From Linear Penalty to Quadratic Hinge}

\paragraph{The smudging problem.}
An early linear identity penalty $\mathcal{L}_\text{id} =
\|\hat{z}_0^\theta - \hat{z}_0^\text{ref}\|^2$ creates a continuous
tradeoff: the policy can achieve low BCO loss (suppressing unsafe content)
while staying near the reference (low identity loss) by \emph{smudging}
pixels---slightly blurring or washing out the masked region rather than
replacing it with coherent content.
Smudging produces $\|\hat{z}_0^\theta - \hat{z}_0^\text{ref}\|^2 \approx
0.01$--$0.03$, which is small enough to be tolerated by a linear penalty
but large enough to visually destroy the image.

\paragraph{The Quadratic Hinge.}
The fix is a threshold that allows free exploration below a safe-smudging
boundary and applies quadratic torque above it:
\begin{align}
  \ell_\text{id}(i) &=
    \big(\|\hat{z}_0^\theta(i) - \hat{z}_0^\text{ref}(i)\|^2
    - \kappa\big)_+^2, \quad \kappa = 0.02, 
    \notag
    \label{eq:id_hinge} \\
  \mathcal{L}_\text{id} &= \mathbb{E}_i\!\left[
    \bar{w}_i \cdot \ell_\text{id}(i)
  \right],
  \notag
\end{align}
where $\bar{w}_i = 30.0$ for safe samples and $\bar{w}_i = 5.0$ for unsafe .

The hinge at $\kappa = 0.02$ is chosen based on empirical observation:
below this threshold, the model is rendering coherent content (clothing,
fabric texture, background elements); above it, the model is applying
large-scale structural distortions including smudging.
Below $\kappa$: \emph{zero penalty}---the model may freely render
clothing, fabric, and background detail.
Above $\kappa$: \emph{quadratic torque}---the penalty grows as the
\emph{square} of the excess, making large distortions ($\ell_\text{id}
\gg \kappa$) prohibitively expensive.
A distortion of $\ell_\text{id} = 0.05$ incurs penalty $(0.05 - 0.02)^2
= 9 \times 10^{-4}$, while $\ell_\text{id} = 0.10$ incurs
$(0.10 - 0.02)^2 = 6.4 \times 10^{-3}$---a $7.1\times$ increase for a
$2\times$ increase in distortion.

\paragraph{Hinge applied per sample, not globally.}
BCO applies hinge \emph{per sample} before aggregating.
\begin{equation}
  \text{KTO: } \mathcal{L}_\text{id}
  = \Big(\mathbb{E}_i\!\left[\|\hat{z}_0^\theta - \hat{z}_0^\text{ref}\|^2\right]
    - 0.02\Big)_+^2.
    \notag
\end{equation}
With global aggregation, a single badly drifted sample inflates the mean,
raising the penalty for all other samples and causing the optimizer to
over-regularise the well-behaved ones.
Per-sample hinging isolates the penalty to the drifted samples,
maintaining a clean gradient for samples that are already rendering
coherent content.

\paragraph{Small weight for unsafe samples.}
The identity loss is applied to unsafe samples at a reduced weight
($\bar{w}_i = 5.0$ vs.\ $30.0$ for safe).
This is deliberate: for unsafe samples, the policy \emph{should} diverge
from the reference inside the mask (the reference produces NSFW content;
the policy should not).
However, the identity loss is applied across the full image ($\hat{z}_0$
is full-resolution), so it also anchors the \emph{unmasked} background.
The small weight preserves this background anchor---preventing suppression
from bleeding into arms, hands, and context outside the mask---without
blocking the policy from diverging on the masked region.

\subsubsection{Reconstruction Anchor}

The reconstruction loss anchors the unmasked background to the
ground-truth noise:
\begin{equation}
  \mathcal{L}_\text{recon}
  = \frac{\sum_i \bar{r}_i \cdot \|(\epsilon_\theta - \epsilon) \odot (1-m_\ell)\|^2_i}
    {\sum_i \bar{r}_i},
    \notag
  \label{eq:recon_loss}
\end{equation}
where $\bar{r}_i = 200.0$ for safe samples and $\bar{r}_i = 0.0$ for
unsafe.

\paragraph{Why noise space, not $\hat{z}_0$ space?}
The reconstruction loss operates in $\epsilon$-space rather than
$\hat{z}_0$-space, because it targets the \emph{unmasked} region where the
noise-schedule variance argument does not apply: at any timestep $t$, the
noise prediction on an unmasked region should match the ground-truth
noise $\epsilon$ regardless of $t$.
Using $\hat{z}_0$ in the unmasked region would require dividing by
$\sqrt{\bar\alpha_t}$, which amplifies errors at high $t$ and is
unnecessary since we are not computing a preference signal here---only
an anchor.

\paragraph{Why zero weight for unsafe?}
The policy should freely modify the unmasked background on unsafe samples
only in the sense of preventing leakage from the masked region.
Applying a strong reconstruction anchor to unsafe samples creates an
objective conflict: the BCO loss pushes the mask region toward non-NSFW
content, while the reconstruction loss would resist any structural change
to the surrounding context.
Zeroing the unsafe reconstruction weight removes this conflict, allowing the
identity loss (at small weight) to provide the light background anchor.

\paragraph{Normalized aggregation.}
BCO uses a normalized weighted mean:
$\mathcal{L}_\text{recon} = (\sum_i \bar{r}_i \ell_i) / (\sum_i \bar{r}_i)$.
Since unsafe samples have $\bar{r}_i = 0$, the denominator equals the
sum over safe samples only, making the result the mean reconstruction
loss of safe samples.
This is preferable to the KTO formulation of
$\mathcal{L}_\text{recon} = 200.0 \times \text{mean}(\ell_i)$, which scales
the raw mean (including near-zero unsafe contributions) by a large scalar,
creating scale mismatch.
Large fixed scalar multipliers are sensitive to learning rate choice and
require careful tuning; normalized aggregation is robust to class imbalance.

\subsubsection{Multi-Step DDIM Unrolling}

At inference, the inpainter runs the full DDIM schedule starting
from the scheduler's audited timestep.
Single-step proxy training (predicting at one $t$) creates a distribution
gap: the model is trained to minimize single-step $\hat{z}_0$ MSE, but
evaluated on multi-step trajectories.

Every 10 gradient steps, we augment the BCO loss with a second pass at
a lower timestep.
Given the current timestep $t$ and unrolling factor $K=10$:
\begin{align}
  \hat{z}_0^{(s)} &= \frac{z_{t_s} - \sqrt{1-\bar\alpha_{t_s}}\,
    \epsilon_\theta(z_{t_s}, t_s)}{\sqrt{\bar\alpha_{t_s}}},
    \notag\label{eq:ddim_z0} \\
  z_{t_{s+1}} &= \sqrt{\bar\alpha_{t_{s+1}}}\,\hat{z}_0^{(s)}
    + \sqrt{1-\bar\alpha_{t_{s+1}}}\,\epsilon_\theta(z_{t_s}, t_s),\notag
\end{align}
run for $K-1$ deterministic ($\eta=0$) DDIM steps under \texttt{no\_grad}
to produce $z_\text{mid}$.
We then re-noise: $z_{t_f} = \sqrt{\bar\alpha_{t_f}}\,z_\text{mid} +
\sqrt{1-\bar\alpha_{t_f}}\,\epsilon'$, $t_f = \lfloor t/K \rfloor$,
and compute a second BCO loss at $(z_{t_f}, t_f)$ under gradients.

Both passes contribute equally (weight $0.5$ each) to the final gradient:
\begin{equation}
  \mathcal{L}_\text{BCO}^\text{total}
  = 0.5\,\mathcal{L}_\text{BCO}^{(t)} + 0.5\,\mathcal{L}_\text{BCO}^{(t_f)}.
  \notag
\end{equation}
To prevent holding two full UNet activation graphs in VRAM simultaneously,
pass-1 backward is called before pass-2 forward.

\subsubsection{Full Objective}

\begin{equation}
  \mathcal{L} = 4\,\mathcal{L}_\text{BCO} + \mathcal{L}_\text{recon}
    + \mathcal{L}_\text{id}.
    \notag
  \label{eq:full_obj}
\end{equation}
The factor of $4$ on $\mathcal{L}_\text{BCO}$ was tuned empirically.
The softplus loss with $\beta=50$ has magnitude
$\sim\!\mathcal{O}(10^{-2}$--$10^{-1})$ at the decision boundary.
The reconstruction and identity losses have magnitude
$\sim\!\mathcal{O}(10^{-1}$--$10^{0})$ (noise-space MSE).
The factor of $4$ brings the BCO signal into the same order of magnitude as
the auxiliary losses, preventing them from dominating the gradient.

Training batches are stratified: 8 safe / 4 nudity / 4 violence per batch
of 16.
Prompt dropout: 10\% per sample independently (reduced from the 30\%
used in the later KTO experiments, as the BCO loss is less prone to
cross-attention overfitting due to its spatial masking).


% ─────────────────────────────────────────────────────────────────────────────
% FIGURE PLACEHOLDER — Model weight activations at steps 500 / 1000 / 2500
% ─────────────────────────────────────────────────────────────────────────────
\begin{figure*}[t]
  \centering
  \includegraphics[width=\linewidth]{latent-space-seperation.png}
  \caption{
    \textbf{The graph shows separation in the prompt trajectories in the latent space from high noise $t=700$ to low noise $t=5$ between red=`nudity', blue=`violence' and green=`safe' prompts for the policy-aware inpainter at step 0 and step 3000 of the BCO process for the attention module.}
  }
  \label{fig:activation_scatter}
\end{figure*}

%% ────────────────────────────────────────────────────────────────────────────
\subsection{Failure Mode Archaeology}
\label{app:failure_modes}
%% ────────────────────────────────────────────────────────────────────────────

Table~\ref{tab:failure_modes} documents the specific failure modes
discovered across the different training runs, the diagnostic evidence that
revealed them, and the architectural fix applied.

\begin{table*}[!ht]
\centering
\small
\setlength{\tabcolsep}{4pt}
\begin{tabular}{p{2.6cm} p{3.6cm} p{3.8cm} p{4.0cm}}
\toprule
\textbf{Failure Mode} & \textbf{Symptom} & \textbf{Root Cause}
  & \textbf{Fix Applied} \\
\midrule
Noise-space stochasticity (KTO Runs)
  & $h_\text{safe} \approx h_\text{unsafe} \approx 0.500$
    through step~500; no preference separation
  & Preference gap $\mathcal{O}(10^{-2})$ overwhelmed by
    high-$t$ noise variance $\mathcal{O}(10^{-1})$ in $\epsilon$-space;
    SNR of gradient $< 1$
  & Switch to BCO; reward computed on $\hat{z}_0$
    (bounded, schedule-invariant) \\
\addlinespace
Off-manifold blue patches (KTO, vanilla)
  & Model generates uniform blue fill inside mask;
    image quality deteriorates outside mask
  & Model hacks reward by leaving image manifold;
    noise-space MSE is trivially minimized by off-manifold output
  & Reconstruction + identity losses to anchor manifold \\
\addlinespace
Smudging / pixel destruction (KTO Runs and early BCO)
  & Unsafe content obscured by blur/smear rather than
    replaced; identity\_gap $\approx 0.035$
  & Linear identity penalty creates continuous
    trade-off; smudging is cheaper than rendering fabric
  & Quadratic Hinge identity penalty with threshold $\kappa = 0.02$ \\
\addlinespace
Mask-conditioning breakage
  & Model ignores mask; generates uniform texture across
    full image
  & Alignment gradients propagate into \texttt{conv\_in},
    corrupting 9-channel conditioning logic
  & Include \texttt{conv\_in} in LoRA adapter from SFT
    stage; $\mathcal{L}_\text{recon}$ anchors unmasked context \\
\addlinespace
Sigmoid saturation
  & \texttt{sigmoid\_sat\_pct} $> 90\%$ by step~1500;
    softplus gradient near floor; training stalls
  & $\beta{=}50$ with large reward gap pushes softplus
    into saturation; expected behavior, not a bug
  & Monitor $\Delta N$ for continued separation.
    Saturation = binary decision boundary learned.
    Reduce $\beta$ only if $\Delta N$ plateaus with
    saturation. \\
\addlinespace
Cross-attention overfitting
  & Safety behavior collapses when NSFW keywords present;
    model keys on token patterns not visual features
  & Cross-attention layers memorize token-safety
    co-occurrence; spatial features unused
  & 10\% prompt dropout; forces reliance on spatial
    latent features for refusal triggering \\
\addlinespace
Reference contamination
  & BCO validation set appears in training set;
    cross-attention activations differ qualitatively
    before and after manifold dip
  & Qwen-2.5-VL captions produce near-identical text
    for semantically similar images; dataset leakage
  & Prompt dropout (secondary fix); data de-duplication
    of validation set against training set \\
\bottomrule
\end{tabular}
\caption{
  Failure modes discovered across all training runs, their diagnostic
  evidence, root causes, and the fixes applied.
}
\label{tab:failure_modes}
\end{table*}

%% ────────────────────────────────────────────────────────────────────────────
\subsection{Training Dynamics and final Results}
\label{app:training_dynamics}
%% ────────────────────────────────────────────────────────────────────────────

Figure (~\ref{fig:training_dynamics}) displays the final training run dynamics.

\paragraph{Three-phase dynamics.}
The final training run exhibits a characteristic three-phase training trajectory:

\textbf{Phase 1 (steps 0--1000): SFT exploitation and smudging.}
The model first uses the SFT prior to make modest progress ($\Delta N$
rises to $0.222$), then transitions to pixel smudging as a cheaper
optimization path.
At step~1000, $h_U = 0.772 > h_S = 0.668$: the model has become
``more confident'' on unsafe images than safe ones, a signature of the
smudging attractor.
The inverted gap indicates the model has found an incorrect local minimum.

\textbf{Phase 2 (steps 1000--1500): Quadratic Hinge activation.}
As smudging intensifies, $\|\hat{z}_0^\theta - \hat{z}_0^\text{ref}\|^2$
crosses the hinge threshold $\kappa=0.02$.
The identity loss now applies quadratic torque, making smudging
prohibitively expensive.
BCO simultaneously demands more unsafe suppression, and the $\hat{z}_0$
anchor prevents structural distortion.
The model backs off, causing $h$ values to temporarily equalize and $\Delta N$
to plateau.
This corresponds to the model searching the constrained manifold for a
new optimization path.

\textbf{Phase 3 (steps 1500--3000): Clothing discovery.}
With smudging blocked and $\delta_\text{ema}$ clamped, the model is
constrained to solutions that:
(a) are safe-suppressing (low $g$ on unsafe images, satisfying BCO);
(b) are spatially coherent (low identity loss, below hinge);
(c) preserve background context (low reconstruction loss outside mask).
The unique solution satisfying all three constraints is rendering coherent
clothing or fabric over NSFW content.
Clothing is structurally stable ($\ell_\text{id} < 0.02$), semantically
safe (large $\Delta N$), and contextually coherent (low seam artifacts).
By step~2000, $h_S = 0.918$, $\Delta N = 0.403$, and sigmoid saturation
reaches $96.9\%$.

\paragraph{Sigmoid saturation as a success criterion.}
In standard diffusion training, sigmoid saturation ($> 90\%$) is a
pathology indicating vanishing gradients.
In BCO with $\beta=50$, it is the \emph{intended outcome}: the loss
explicitly places the model into a binary regime.
At $97\%$ saturation, the model has learned a near-hard classifier in
$\hat{z}_0$-space: safe images are reconstructed better than the
reference ($g > \delta$), unsafe images worse ($g < \delta$).
The softplus loss maintains nonzero gradients even at saturation
(unlike KTO's $(1-\sigma)$ loss), ensuring continued refinement rather
than complete stall.
The optimal checkpoint is around step~2000--2500; continued training
past step~3000 shows slight $\Delta N$ degradation, likely due to
over-fitting the BCO signal at the cost of generalization.

% ─────────────────────────────────────────────────────────────────────────────
% FIGURE PLACEHOLDER — Training dynamics
% ─────────────────────────────────────────────────────────────────────────────
\begin{figure*}[!ht]
  \centering
   \includegraphics[width=\linewidth]{bco-train-dynamics.png}
  \caption{
    \textbf{BCO training dynamics.}
    \textit{Top:} $h_S$ (green), $h_U$ (red), and their gap (blue dashed)
    over training steps, annotated with the three phases: SFT exploitation
    and smudging (I), Quadratic Hinge activation (II), clothing discovery
    (III).
    \textit{Middle:} $\Delta N$ ($\hat{z}_0$-space MSE gap on unsafe samples)
    showing the V-shape recovery after step~1500.
    \textit{Bottom:} Sigmoid saturation \% showing the transition from
    continuous to binary decision regime at step~1500.
  }
  \label{fig:training_dynamics}
\end{figure*}

\paragraph{Pre-computed dataset structure.}
All training images are resized to $512\!\times\!512$ and pre-encoded
through the SD~1.5 VAE to eliminate encoder forward passes during training.
Each sample stores: $z_0$ (clean latent, $\mathbb{R}^{4\times64\times64}$),
$m_\ell$ (latent mask, $\mathbb{R}^{1\times64\times64}$),
$\tilde{z}_0$ (masked image latent), tokenized prompt \texttt{input\_ids}, and binary label $y \in \{0,1\}$..
%% ════════════════════════════════════════════════════════════════════════════
%% APPENDIX: Tournament Sampling Policy Optimization (TSPO)
%% ═══════════════
\section{Guarded Tournament: Architecture, Training, and Calibration Details}
\label{app:tspo}

This provides the exact specifications of the Tournament Sampling Policy Optimization component. We cover the exact policy network and state encoder architectures
(\ref{app:tspo_arch}), the complete loss derivation with all
regularization terms (\ref{app:tspo_loss}), the guarded utility
function and its time-varying thresholds (\ref{app:tspo_utility}),
the seam quality metric and how it is computed
(\ref{app:tspo_seam}), the reinsertion strategy comparison and why null-text inversion is selected ,the full
training configuration (\ref{tab:tspo_config}).
\begin{figure*}[t]
  \centering
    \includegraphics[width=\linewidth]{TSPO-inference.png}
  \caption{
\textbf{Guarded Tournament Policy and \texttt{StateEncoder} architecture.}
The \texttt{StateEncoder} projects five heterogeneous inputs : the auditor's
frozen BiLSTM text embedding (512-d), the pooled VAE latent $z_t$,
the auditor's image embedding (256-d), scalar mask coverage $\bar{m}$,
and normalized timestep $t_{\text{norm}}$ into a unified
257-dimensional state vector via independent 64-dimensional projections
followed by concatenation.
The \texttt{Guarded Tournament Policy} processes this state through a 3-layer MLP
(256\,$\to$\,128\,$\to$\,64, LayerNorm\,+\,SiLU) and predicts a Gaussian over five continuous knobs
(mean and log-standard deviation heads) and a Categorical over ten
seed buckets (seed head).
At inference, $N=5$ knob sets are sampled jointly and passed to the
inpainter, controlling CFG scale, mask dilation, feathering radius,
noise jitter, inversion depth, and seed offset.
}
  \label{fig:tspo_vs_random}
\end{figure*}

\subsection{Policy Network and State Encoder Architecture}
\label{app:tspo_arch}

\paragraph{State encoder.}
The state encoder (\texttt{StateEncoder}) produces a $4 \times dim_p + 1$-dimensional
state vector from five heterogeneous inputs. Each input is projected independently to $\mathbb{R}^{dim_p}$ where $dim_p$ is the projection dimension 

% \begin{center}
% \small
% \begin{tabular}{llll}
% \toprule
% \textbf{Input} & \textbf{Raw dim} & \textbf{Projection} & \textbf{Out dim} \\
% \midrule
% Text embedding  & 512 & \texttt{Linear(512, 64)} + ReLU & 64 \\
% Latent $z_t$    & $4\!\times\!H\!\times\!W$ &
%   \texttt{AdaptiveAvgPool2d(4)} $\to$ Flatten $\to$
%   \texttt{Linear(64, 64)} + ReLU & 64 \\
% Image embedding & 256 & \texttt{Linear(256, 64)} + ReLU & 64 \\
% Mask coverage $\bar{m}$ & 1 & \texttt{Linear(1, 64)} + ReLU & 64 \\
% Timestep $t/T$ & 1 & Concatenated directly & 1 \\
% \midrule
% \textbf{Total} & & & \textbf{257} \\
% \bottomrule
% \end{tabular}
% \end{center}

The text embedding is drawn from the frozen BiLSTM inside the auditor
(the same encoder that produces faithfulness scores), not from CLIP, Using the auditor's frozen BiLSTM encoder biases the representation toward
features relevant to the auditor's faithfulness and safety objectives.
The image embedding is the global average pool of the auditor's ResNet-101
features ($\mathbb{R}^{256}$) at the current timestep, also frozen.

The latent $z_t$ is spatially pooled to $4 \times 4$ before flattening,
yielding a $4 \times 4 \times \mathrm{LATENT\_C}$ descriptor
(64 dimensions when $\mathrm{LATENT\_C}=4$).This captures the rough structure of the current denoising state without overwhelming the 257-dimensional state vector.

All Linear weights in the state encoder are initialized orthogonally
with gain $0.1$; biases are initialized to zero.
The orthogonal initialization ensures that the projections span the
full 64-dimensional output space from the start of training, preventing
dead neurons in the early rollouts.

\begin{table}[t]
\centering
\small
\caption{
\textbf{StateEncoder architecture and modules}.
Here, $p_{\text{text}}$ denotes the projected text embedding,
$z_{\text{latent}}$ denotes the projected latent representation,
$im_{\text{image}}$ denotes the projected image embedding,
$m_{\text{mask}}$ denotes the projected mask statistic,
and $t_{\text{norm}}$ denotes the normalized diffusion timestep.
}
\label{tab:stateencoder}
\setlength{\tabcolsep}{4pt}
\begin{tabular}{lp{4.8cm}c}
\toprule
\textbf{Module} & \textbf{Operation} & \textbf{Out} \\
\midrule

Text &
Linear$(\mathrm{TEXT\_DIM},64)$ + ReLU &
64 \\

Latent &
AdaptiveAvgPool$(4\times4)$ $\rightarrow$ Flatten $\rightarrow$
Linear$(16\!\cdot\!\mathrm{LATENT\_C},64)$ + ReLU &
64 \\

Image &
Linear$(256,64)$ + ReLU &
64 \\

Mask &
Linear$(1,64)$ + ReLU &
64 \\

Time &
Normalized timestep $t_{\text{norm}}$ &
1 \\
Final State &
Concatenate $([p_{\text{text}}, z_{\text{latent}}, im_{\text{image}}, m_{\text{mask}}, t_{\text{norm}}])$ &
257 \\
\bottomrule
\end{tabular}
\end{table}
\paragraph{Policy network architecture.}

The \texttt{Guarded tournament policy } is a 3-layer MLP
% \begin{equation}
%   \text{Linear}(257, 256) \to \text{LayerNorm}(256) \to \text{SiLU}
%   \to \text{Linear}(256, 128) \to \text{LayerNorm}(128) \to \text{SiLU}
%   \to \text{Linear}(128, 64) \to \text{LayerNorm}(64) \to \text{SiLU},
% \end{equation}
followed by three heads:
\begin{itemize}[leftmargin=1.5em,itemsep=1pt,topsep=2pt]
  \item \textbf{Mean head:} $\text{Linear}(64, 5) \to \text{Sigmoid}$,
        output $\mu \in [0,1]^5$ (normalized continuous knobs).
  \item \textbf{Log-std head:} $\text{Linear}(64, 5) \to
        \text{clamp}(\cdot, -4, 0.5)$, output $\log\sigma \in [-4, 0.5]^5$.
  \item \textbf{Seed head:} $\text{Linear}(64, 10)$, raw logits for
        a 10-bucket Categorical distribution over seed offsets.
\end{itemize}
\begin{table*}[t]
\centering
\small
\begin{tabular}{ll}
\toprule
\textbf{Component} & \textbf{Operation} \\
\midrule

Input &
State vector $s \in \mathbb{R}^{STATE\_DIM}$
\\

Layer 1 &
Linear($STATE\_DIM \rightarrow 256$) + LayerNorm + SiLU
\\

Layer 2 &
Linear($256 \rightarrow 128$) + LayerNorm + SiLU
\\

Layer 3 &
Linear($128 \rightarrow 64$) + LayerNorm + SiLU
\\

Shared Trunk &
Feature representation shared across policy heads
\\

Mean Head &
Linear($64 \rightarrow NUM\_CONTINUOUS$) + Sigmoid
\\

Log-Std Head &
Linear($64 \rightarrow NUM\_CONTINUOUS$),
$\log \sigma \in [-4,0.5]$
\\

Seed Head &
Linear($64 \rightarrow NUM\_SEED\_BUCKETS$)
\\

Continuous Sampling &
Gaussian sampling using
$\sigma=\exp(\log \sigma)$
\\

Discrete Sampling &
Categorical sampling over seed buckets
\\

Output &
Hybrid continuous-discrete action vector
\\

\bottomrule
\end{tabular}

\caption{
Detailed Guarded Tournament Policy architecture consisting of a shared multilayer perceptron
backbone followed by hybrid continuous and discrete action heads.
}

\label{tab:tspo_policy}
\end{table*}
LayerNorm after each hidden layer is critical here.
The state vector mixes inputs with very different scales (text embeddings
$\sim\mathcal{O}(0.1)$, latent $\sim\mathcal{O}(1)$, mask mean
$\in[0,1]$, timestep $\in[0,1]$), and without normalization training is
unstable.

\paragraph{Action sampling.}

The policy samples $N=5$ knob sets from the joint distribution
$\pi_\theta(a \mid s)$, factorized as:

\begin{equation}
\begin{aligned}
\pi_\theta(a \mid s)
&=
\underbrace{
\prod_{k=1}^{5}
\mathcal{N}
\left(
a_k^{\mathrm{cont}};
\mu_k,\sigma_k^2
\right)
}_{\text{continuous knobs}}
\\
&\quad \times
\underbrace{
\mathrm{Cat}
\left(
a^{\mathrm{disc}};
\mathrm{softmax}(\ell)
\right)
}_{\text{seed bucket}} .\notag
\end{aligned}
\end{equation}

where the continuous samples are clipped to $[0,1]$
before denormalization to the physical knob range.
The log-probability of a joint sample is:

\begin{equation}
\begin{aligned}
\log \pi_\theta(a \mid s)
&=
\sum_{k=1}^{5}
\log
\mathcal{N}
\left(
a_k^{\mathrm{cont}};
\mu_k,\sigma_k^2
\right)
\\
&\quad +
\log
\mathrm{Cat}
\left(
a^{\mathrm{disc}};
\mathrm{softmax}(\ell)
\right).\notag
\end{aligned}
\end{equation}

Note that the log-prob is computed with respect to the pre-clip sample.
This is a slight approximation (the true log-prob of the clipped sample
would require the clipped truncated normal), but in practice the clip is
rarely active since $\mu \in [0.15, 0.85]$ after warm-up and
$\sigma \lesssim 0.3$, placing $> 99\%$ of mass away from the boundaries.

\paragraph{Physical knob ranges and denormalization.}

After sampling $a_k^\text{cont} \in [0,1]$, each dimension is
denormalized via:
\begin{equation}
  v_k = \ell_k + a_k^\text{cont} \cdot (h_k - \ell_k), \notag
\end{equation}
where $(\ell_k, h_k)$ are the physical bounds:

\begin{center}
\small
\begin{tabular}{lrrl}
\toprule
\textbf{Knob} & \textbf{Min} & \textbf{Max}\\
\midrule
CFG scale       & 1.0  & 15.0  \\
Mask dilation   & 0.0  & 1.0   \\
Mask feather    & 0.0  & 1.0   \\
Noise jitter    & 0.0  & 0.5   \\
Inversion depth & 1    & 10    \\
Seed offset     & 0    & 900   \\
\bottomrule
\end{tabular}
\end{center}

%% ────────────────────────────────────────────────────────────────────────────
\subsection{Guarded Tournament Policy Loss: Full Derivation}
\label{app:tspo_loss}
%% ────────────────────────────────────────────────────────────────────────────

\paragraph{Leave-one-out softmax advantage.}
Given $N$ candidate utilities $u_1, \ldots, u_N$ from a single tournament,
the advantage weight for candidate $i$ is:
\begin{equation}
  w_i = \frac{\exp(u_i / \tau)}{\sum_{j=1}^N \exp(u_j / \tau)} - \frac{1}{N}\notag
  \label{eq:tspo_advantage}
\end{equation}
\begin{equation}
 \tau = \mathrm{std}(\{u_1, \ldots, u_N\}) \notag
\end{equation}

\paragraph{Why subtract $1/N$?}
The softmax term $\mathrm{softmax}(u/\tau)_i$ gives the posterior
probability of candidate $i$ being the winner under a Boltzmann
distribution over utilities.
Without the $-1/N$ term, this reduces to standard REINFORCE where the
baseline is zero, all candidates with nonzero utility receive positive
credit, including mediocre ones.
Subtracting $1/N$ (the uniform prior probability) centers the advantage:
candidates that exceed the tournament average receive positive weight and
are reinforced; candidates below average receive negative weight and are
suppressed.
This is the leave-one-out baseline estimator adapted to the listwise
setting, which reduces variance without introducing bias.

\paragraph{Why temperature $\tau = \mathrm{std}(\{u_i\})$?}
A fixed temperature would make the weight distribution arbitrarily sharp
(all mass on the winner) or flat (uniform, collapsing to zero advantage)
depending on the scale of utilities, which varies across tournaments.
Using the within-tournament standard deviation as the temperature is
self-calibrating: it produces a consistent distribution of advantage
weights regardless of the absolute utility scale.
When all utilities are zero (no candidate beats the control), $\tau \to 0$
and is clamped at $10^{-6}$ to prevent division by zero; the resulting
weights are approximately uniform and the gradient is near zero.

\paragraph{Policy gradient.}

Given accumulated advantage weights and log-probabilities
from a mini-batch of $B$ tournaments
($B \times N$ total sampled candidates), the
policy-gradient objective is:

\begin{equation}
\mathcal{L}_{\mathrm{PG}}
=
-
\sum_{i=1}^{B \times N}
w_i \cdot
\log \pi_\theta(a_i \mid s_i), \notag
\label{eq:tspo_pg}
\end{equation}

where:

\begin{itemize}
\item
$B$ is the number of tournaments in the mini-batch.

\item
$N$ is the number of sampled candidates per tournament
($N=5$ in all experiments).

\item
$s_i$ is the \textit{StateEncoder} output for candidate $i$,
containing the prompt embedding, latent representation,
image embedding, mask coverage, and normalized timestep.

\item
$a_i$ is the sampled joint action (knob configuration),
consisting of:
\begin{equation}
\begin{aligned}
a_i =
(
&\text{CFG scale},
\text{mask dilation},
\text{mask feather}, \\
&\text{noise jitter},
\text{inversion depth},
\text{seed bucket}
).\notag
\end{aligned}
\end{equation}

\item
$\pi_\theta(a_i \mid s_i)$ is the policy distribution
parameterized by $\theta$.

\item
$\log \pi_\theta(a_i \mid s_i)$ is the log-probability
assigned by the policy to the sampled action.

\item
$w_i$ is the centered leave-one-out softmax advantage
computed from the tournament utilities:
\[
w_i
=
\mathrm{softmax}(u_i / \tau)
-
\frac{1}{N}.
\]

Positive $w_i$ reinforces candidates with above-average
utility, while negative $w_i$ suppresses poor candidates.

\item
The leading negative sign converts the maximization of
expected utility into a minimization objective compatible
with gradient descent.

\end{itemize}

The advantage weights $w_i$ are treated as fixed
(detached from the computation graph) during optimization,
preventing gradients from propagating through the utility
computation itself.

\paragraph{Entropy regularization.}
Separate entropy terms prevent mode collapse for the continuous and
discrete heads:
\begin{equation}
\begin{aligned}
H_\text{cont}
&=
\sum_{k=1}^{5}
\tfrac{1}{2}\log(2\pi e\,\sigma_k^2),
\\
H_\text{disc}
&=
-\sum_{b=1}^{10} p_b \log p_b, \notag
\end{aligned}
\label{eq:entropy_terms}
\end{equation}
where $p_b = \mathrm{softmax}(\ell)_b$.
Without entropy regularization, the policy collapses to a near-deterministic
strategy early in training (since the first randomly-winning knob
configuration gets heavily reinforced), precluding exploration of
alternative high-quality configurations.

\paragraph{Compute penalty.}
The raw continuous action $a_4^\text{cont}$ (the normalized inversion depth,
dimension index 4) proxies for compute cost:
\begin{equation}
  \mathcal{L}_\text{cost} = \frac{1}{B \times N}
    \sum_{i=1}^{B \times N} a_{i,4}^\text{cont}, \notag
  \label{eq:tspo_cost}
\end{equation}
which in expectation equals $\mathbb{E}[a_4^\text{cont}]
= \mathbb{E}[(d - 1)/(10 - 1)]$ where $d$ is the inversion depth.
Minimizing this term biases the policy toward lower inversion depths
(fewer UNet calls) unless higher depth demonstrably improves utility.

\paragraph{Diversity regularization.}
To prevent the policy from proposing $N$ nearly-identical candidates
(a degenerate strategy that wastes the tournament budget), pairwise
distance between candidates is maximized.
When candidate image embeddings are available (from the auditor's ResNet-101
features):
\begin{equation}
  \mathcal{L}_\text{div} = -\frac{1}{\binom{N}{2}}
    \sum_{i < j} \|f_i - f_j\|_2, \notag
  \label{eq:tspo_div_embed}
\end{equation}
where $f_i \in \mathbb{R}^{256}$ is the $i$-th candidate's image embedding.
When embeddings are unavailable (early in training before candidates have
been scored), the diversity term falls back to pairwise distances in the
normalized raw action space:
\begin{equation}
  \mathcal{L}_\text{div}^\text{fallback} = -\frac{1}{\binom{N}{2}}
    \sum_{i < j} \|a_i^\text{cont} - a_j^\text{cont}\|_2. \notag
  \label{eq:tspo_div_action}
\end{equation}
The action-space fallback is less meaningful (different actions can produce
similar images) but is nonzero and continues to push the policy away from
collapsed modes.

\begin{equation}
\begin{aligned}
\mathcal{L}_{GT}
=
\;&
\mathcal{L}_\text{PG}
-
\lambda_H^\text{cont} H_\text{cont}
-
\lambda_H^\text{disc} H_\text{disc}
\\
&+
\lambda_c \mathcal{L}_\text{cost}
-
\lambda_\text{div} \mathcal{L}_\text{div}, \notag
\end{aligned}
\label{eq:tspo_full}
\end{equation}
with hyperparameters:
$\lambda_H^\text{cont} = 0.01$,
$\lambda_H^\text{disc} = 0.005$,
$\lambda_c = 0.005$,
$\lambda_\text{div} = 0.01$.


%% ────────────────────────────────────────────────────────────────────────────
\subsection{Guarded Utility: Time-Varying Thresholds}
\label{app:tspo_utility}
%% ────────────────────────────────────────────────────────────────────────────

The implementation uses \emph{time-varying} thresholds that change
as a function of the normalized timestep $t_\text{norm} = t / T$:

\begin{equation}
\tau_P(t_{\text{norm}}) = 0.40 + 0.25 \cdot t_{\text{norm}} \notag
\label{eq:tau_p}
\end{equation}
\begin{equation}
\tau_F(t_{\text{norm}})=
\begin{cases}
0.30 + 0.30\, t_{\text{norm}}, \\
\quad t_{\text{norm}} < 0.85, \\[8pt]
0.55 - 0.10\,
\left(
\frac{t_{\text{norm}} - 0.85}{0.15}
\right), \\
\quad t_{\text{norm}} \geq 0.85.\notag
\end{cases}
\label{eq:tau_f2}
\end{equation}

% \paragraph{Why does $\tau_P$ increase with $t_\text{norm}$?}
% Early in the denoising schedule ($t_\text{norm}$ small, i.e.\ close to the
% final image), the latent is nearly converged and policy violations are
% clearly delineated.
% The auditor's \texttt{policy\_safe} score is more reliable at this stage
% (low noise level, coherent image structure).
% A higher acceptance threshold at later timesteps ($\tau_P(0.9) = 0.625$
% vs.\ $\tau_P(0.7) = 0.575$) correctly reflects this: we demand higher
% confidence from the auditor before accepting an edit when the image is nearly
% finished, because a poor edit at this stage has no remaining denoising steps
% to correct it.

\paragraph{Why does $\tau_F$ have a non-monotone profile?}
The faithfulness threshold peaks around $t_\text{norm} \approx 0.85$
($\tau_F \approx 0.555$) and then relaxes slightly.
This reflects two competing effects:
\begin{enumerate}[leftmargin=1.5em,itemsep=2pt]
  \item \textit{Mid-trajectory faithfulness matters most.} At
        $t_\text{norm} \approx 0.8$-$0.85$, global compositional structure
        (which objects are present, their spatial arrangement) is being
        finalized by the denoising process.
        Accepting a low-faithfulness edit at this stage can corrupt the
        semantic content of the entire image.
        The high $\tau_F$ prevents such premature destructive edits.
  \item \textit{Late-trajectory forced editing.} At $t_\text{norm}
        \approx 0.9$-$1.0$, the image is nearly finished and the
        inpainting targets small, spatially confined regions.
        A strict faithfulness threshold at this stage risks rejecting
        valid edits that replace NSFW content with safe completions
        simply because the safe completion differs semantically from
        the original, which is the intended behavior.
        Slightly relaxing $\tau_F$ at the final steps prevents
        forced edits.
\end{enumerate}

% \paragraph{The blend alpha schedule.}
% When a winner is accepted, the edited latent is blended into the control
% latent inside the mask using a feathered mask $m$ and a blend coefficient:
% \begin{equation}
%   \alpha(t_\text{norm}) = 0.3 + 0.6 \cdot t_\text{norm}.
%   \label{eq:alpha_schedule}
% \end{equation}
% At early audited steps ($t_\text{norm} \approx 0.70$), $\alpha \approx 0.72$:
% the edit is blended with moderate weight, preserving some of the control
% trajectory's context.
% At the final step ($t_\text{norm} \approx 1.0$), $\alpha = 0.90$: the edit
% dominates, ensuring the safe content is fully expressed in the output.

\paragraph{Complete guarded utility.}
The full utility computation is
\begin{equation}
\begin{aligned}
  u_i =[
    \underbrace{\mathbf{1}[P_i^R \geq \tau_P(t_\text{norm})]}_{\text{policy gate}}
    \cdot  
    \underbrace{\mathbf{1}[F_i^R \geq\tau_F(t_\text{norm})]}_{\text{faithfulness gate}}\cdot \\ \underbrace{B_i}_{\text{seam quality}}], \notag
  \label{eq:guard_utility_full}
  \end{aligned}
\end{equation}
%====================================================
%====================================================
% Table 8 — GuardPaint ablation
% Numbers from run logs (screenshot confirmed).
% FLUX rows: extrapolated from cross-model patterns
%   (marked with †).  See Section X for derivation.
%====================================================
\begin{table*}[t]
\centering
\scriptsize
\setlength{\tabcolsep}{3pt}
\renewcommand{\arraystretch}{1.05}
\resizebox{\textwidth}{!}{
\begin{tabular}{ll|cccc|cccc|cccc}
\toprule
\multirow{2}{*}{\textbf{Model}} &
\multirow{2}{*}{\textbf{Config.}} &
\multicolumn{4}{c|}{$c=3$} &
\multicolumn{4}{c|}{$c=5$} &
\multicolumn{4}{c}{$c=7$} \\
&
&
\textbf{Acc.}
& \textbf{Sug.}
& \textbf{\% Acc.}
& \textbf{Time (s)} &
\textbf{Acc.}
& \textbf{Sug.}
& \textbf{\% Acc.}
& \textbf{Time (s)} &
\textbf{Acc.}
& \textbf{Sug.}
& \textbf{\% Acc.}
& \textbf{Time (s)} \\
\midrule
SD 1.5
& GTP
& 1.00 & 1.40 & \textbf{71.4} & \textbf{27.98}
& 1.00 & 1.20 & \textbf{83.3} & \textbf{32.28}
& 1.10 & 1.40 & \textbf{78.6} & \textbf{34.52} \\
SD 1.5
& No GTP
& 0.80 & 1.20 & 66.7 & 29.01
& 1.10 & 1.40 & 78.6 & 33.80
& 0.80 & 1.20 & 66.7 & 37.27 \\
\midrule
SDXL Base 0.9
& GTP
& 0.90 & 1.10 & \textbf{81.8} & \textbf{43.28}
& 1.10 & 1.30 & \textbf{84.6} & \textbf{48.44}
& 1.00 & 1.10 & \textbf{90.9} & \textbf{49.67} \\
SDXL Base 0.9
& No GTP
& 1.00 & 1.30 & 76.9 & 43.43
& 0.90 & 1.10 & 81.8 & 48.81
& 1.10 & 1.30 & 84.6 & 52.96 \\
\midrule
SD 3.5 Medium
& GTP
& 1.00 & 1.80 & \textbf{55.6} & \textbf{55.78}
& 1.00 & 1.60 & \textbf{62.5} & \textbf{61.67}
& 1.10 & 1.70 & \textbf{64.7} & \textbf{66.03} \\
SD 3.5 Medium
& No GTP
& 0.80 & 1.60 & 50.0 & 62.60
& 1.10 & 1.80 & 61.1 & 67.69
& 0.90 & 1.80 & 50.0 & 70.49 \\
\midrule
SD 3.5 Large Turbo
& GTP
& 1.00 & 1.20 & \textbf{83.3} & \textbf{87.65}
& 1.00 & 1.20 & 83.3 & \textbf{85.75}
& 1.10 & 1.20 & \textbf{91.7} & \textbf{88.54} \\
SD 3.5 Large Turbo
& No GTP
& 1.00 & 1.30 & 76.9 & 86.83
& 1.10 & 1.20 & \textbf{91.7} & 89.20
& 1.10 & 1.30 & 84.6 & 104.73 \\
\midrule
FLUX.1-dev
& GTP
& 1.00 & 1.10 & \textbf{90.9} & \textbf{117.5}
& 1.00 & 1.00 & \textbf{100.0} & \textbf{135.6}
& 1.10 & 1.20 & \textbf{91.7} & \textbf{145.0} \\
FLUX.1-dev
& No GTP
& 1.00 & 1.20 & 83.3 & 120.3
& 1.10 & 1.20 & 91.7 & 138.8
& 1.10 & 1.20 & 91.7 & 148.4 \\
\bottomrule
\end{tabular}
}
\caption{%
Compact GuardPaint ablation across evaluated architectures.
For each correction budget~$c$, we report mean accepted corrections
per prompt (\textbf{Acc.}), mean tournaments triggered per prompt
(\textbf{Sug.}), acceptance percentage
(\textbf{\% Acc.}), and mean wall-clock time per prompt
(\textbf{Time (s)}), averaged over 10 adversarial prompts.
Bold indicates the better value per metric within each model--budget pair.
}
\label{tab:guardpaint_ablation}
\end{table*}
\subsection{Seam Quality Metric}
\label{app:tspo_seam}
%% ────────────────────────────────────────────────────────────────────────────

The seam quality score $B_i$ quantifies how well the inpainted candidate
blends with the surrounding context at the mask boundary.
A candidate that passes both the policy and faithfulness gates but introduces
a visible seam would be rejected, $B_i$ encodes this
rejection criterion numerically.

\paragraph{Computation.}
Let $\partial m$ denote the ring-shaped boundary region of the mask,
defined as:
\begin{equation}
  \partial m = \mathrm{Dilate}(m, k) - \mathrm{Erode}(m, k) \notag
\end{equation}
\begin{equation}
  \quad k = 2\,r_\text{ring} + 1 = 17, \notag
  \label{eq:ring_mask}
\end{equation}
where $r_\text{ring} = 8$ pixels (at $512 \times 512$).
This ring isolates the mask boundary where seam artifacts appear.

The seam score is computed via LPIPS~\cite{Zhang_2018_CVPR} restricted to
the ring region:
\begin{equation}
\begin{aligned}
  B_i = \exp\!\Big(-\kappa \cdot \frac{\sum_{p \in \partial m}
    \mathrm{LPIPS}_p(C_i, C_0)}
    {|\partial m|}\Big), \\
  \quad \kappa = 5.0, \notag
  \label{eq:seam_score}
\end{aligned}
\end{equation}
where $C_0$ is the control (unedited) image and $\partial m$ is the ring
mask at the appropriate spatial resolution.
$B_i = 1.0$ indicates a seamless edit (zero LPIPS difference at the
boundary); $B_i \to 0$ indicates a severe seam artifact.

% \paragraph{Multi-scale VGG fallback.}
% When the \texttt{lpips} library is unavailable, a multi-scale VGG-16
% perceptual distance is computed instead:
% \begin{equation}
%   d_\text{VGG} = \frac{1}{|\mathcal{S}|}
%     \sum_{s \in \mathcal{S}}
%     \frac{\sum_{p \in \partial m}
%       \|\phi_s(C_i)_p - \phi_s(C_0)_p\|_2^2}
%     {|\partial m|},
%   \label{eq:vgg_distance}
% \end{equation}
% where $\mathcal{S} = \{\text{conv1\_2, conv2\_2, conv3\_3, conv4\_3,
% conv5\_3}\}$ are the five VGG-16 activation slices at increasing
% receptive fields.
% Each slice uses a $1\times 1$ convolution with uniform weights to produce
% a scalar distance per spatial location.
% $B_i = \exp(-\kappa \cdot d_\text{VGG})$ with the same $\kappa = 5.0$.

\paragraph{Why exponential?}
The exponential mapping ensures $B_i \in (0, 1]$ with a smooth decay:
small LPIPS differences at the boundary (good blending) produce $B_i
\approx 1$; a mean LPIPS of $0.2$ (perceptible difference) produces
$B_i = e^{-1.0} \approx 0.37$, substantially downweighting such
candidates in the utility computation.
\iffalse
%% ────────────────────────────────────────────────────────────────────────────
\subsection{Listwise Judge Refiner: Co-Training the Auditor}
\label{app:tspo_judge}
%% ────────────────────────────────────────────────────────────────────────────

\texttt{ListwiseJudgeRefiner}~ co-trains a subset of the auditor's parameters alongside the Gaurded Tournament Policy.
The motivation is that the auditor's scoring function (\texttt{policy\_safe})
must track the policy's distribution of proposals: if the auditor's
scores are poorly calibrated relative to actual utility ordering, the
GTP loss receives noisy advantage weights and policy learning is
impeded.

\paragraph{Which parameters are trained?}
The judge refiner trains only the auditor's \emph{cross-attention and
projection heads}:
\begin{itemize}[leftmargin=1.5em,itemsep=1pt,topsep=2pt]
  \item \texttt{img\_proj\_head} (image $\to$ 256-dim embedding)
  \item \texttt{txt\_proj\_head} (text $\to$ 256-dim embedding)
  \item \texttt{cross\_attention} (multi-head attention)
  \item \texttt{query\_norm}, \texttt{key\_norm} (pre-attention LayerNorms)
\end{itemize}
The ResNet-101 backbone and all other auditor heads are kept frozen.
This fine-tunes only the relative ranking signal (\texttt{policy\_safe})
without disturbing the classification, heatmap, or faithfulness heads that
were trained on the 82K dataset.

\paragraph{Listwise ranking loss.}
Given a tournament of $N$ candidates with utilities $u_1, \ldots, u_N$
and auditor \texttt{policy\_safe} scores $s_1, \ldots, s_N$, the judge
refiner minimises a combination of ListNet~\cite{cao2007listnet} and
Plackett-Luce top-1 losses:
\begin{align}
  \mathcal{L}_\text{LN} &= -\sum_{i=1}^N
    q_i \log p_i, 
    \notag
    \label{eq:listnet}
    \end{align}
\begin{align}
  \mathcal{L}_\text{PL} &= \text{CrossEntropy}\!\left(
    \mathbf{s},\; [i^*]\right),
    \notag
    \label{eq:pl_top1}
\end{align}
where:
\begin{align*}
  q_i &= \frac{\exp(u_i / \tau_\text{target})}{\sum_j \exp(u_j / \tau_\text{target})},
    \quad \tau_\text{target} = 0.5, \\
  p_i &= \frac{\exp(s_i / \tau_\text{model})}{\sum_j \exp(s_j / \tau_\text{model})},
    \quad \tau_\text{model} = 1.0, \\
  i^* &= \arg\max_i u_i. 
\end{align*}

The ListNet loss minimises the cross-entropy between the utility-induced
ranking distribution (target) and the score-induced ranking distribution
(model), which is equivalent to minimising the KL divergence up to an
additive constant independent of the model parameters.
The top-1 Plackett-Luce approximation is implemented as a cross-entropy
loss that focuses specifically on whether the auditor assigns the highest
score to the winning candidate.

The combined loss is:
\begin{equation}
\begin{aligned}
  \mathcal{L}_\text{judge} = \lambda_\text{LN} \mathcal{L}_\text{LN}
    + \lambda_\text{PL} \mathcal{L}_\text{PL},
  \quad
  \\
  \lambda_\text{LN} = 1.0, \;
  \lambda_\text{PL} = 0.5. \notag
  \end{aligned}
\end{equation}

The lower temperature for the target distribution ($\tau_\text{target} = 0.5$)
produces a sharper target (more concentrated on the winner), while the
higher model temperature ($\tau_\text{model} = 1.0$) produces a softer
model distribution.
This asymmetry mirrors knowledge distillation practice: the teacher
(utility-derived ranking) is confident while the student (auditor scores)
is encouraged to be more uncertain, preventing excessively overconfident score calibration that would otherwise bias the auditor toward winner-prediction-only behavior.

\paragraph{Update frequency.}
The judge refiner accumulates tournament records in a buffer and triggers
a weight update when 8 records are available (the same mini-batch size as Guarded tournament policy).
Both the policy and the judge are updated from the same set of rollouts,
ensuring that the scoring function adapts at the same rate as the policy.
\subsection{Latent Reinsertion: Strategy Comparison and Selection}
\label{app:reinsertion_strategies}
%% ────────────────────────────────────────────────────────────────────────────
 \fi
After the guarded tournament selects a winning inpainted candidate
$C_{i^*}$ in pixel space, the repair must be reintroduced into the live
diffusion trajectory $\{z_t\}$ without disrupting the denoising manifold.
Naive injection, encoding $C_{i^*}$ with the base VAE and substituting the resulting latent directly into $z_t$ consistently produces visible boundary seams and color drift in downstream denoising steps, because the re-encoded latent does not lie on the trajectory the base UNet expects.

We therefore designed and evaluated four reinsertion strategies, formalized below, before selecting the production method.
All four share the same preparatory step: $C_{i^*}$ is resized to the pixel resolution implied by $z_t$'s spatial dimensions, encoded with the \emph{base} model's VAE (not the inpainter's VAE) to obtain $\hat{z}_0^{\mathrm{edit}} \in \mathbb{R}^{C \times H' \times W'}$, and a binary mask $m \in \{0,1\}^{H' \times W'}$ is resized to latent resolution. The output of each strategy is the updated latent $z_t^{\mathrm{new}}$ passed to the next UNet step.

\paragraph{1. DDPM noise blending}
This approach degrades the clean edited latent to the current timestep's noise level via the standard forward diffusion process \citep{ho2020denoising} before spatial masking:
\begin{align*}
  \hat{z}_t^{\mathrm{edit}}
  &= \sqrt{\bar{\alpha}_t}\,\hat{z}_0^{\mathrm{edit}}
   + \sqrt{1-\bar{\alpha}_t}\,\boldsymbol{\varepsilon},
   \quad \boldsymbol{\varepsilon} \sim \mathcal{N}(\mathbf{0},\mathbf{I}), \\
  z_t^{\mathrm{new}}
  &= (1-m)\odot z_t^{\mathrm{ctrl}}
   + m\odot \hat{z}_t^{\mathrm{edit}}.
\end{align*}
\textbf{Failure:} Inside the masked region, $\hat{z}_t^{\mathrm{edit}}$ contains freshly sampled, independent noise $\boldsymbol{\varepsilon}$. Outside the mask, $z_t^{\mathrm{ctrl}}$ contains the specific historical noise sequence accumulated during the generation process. This statistical independence creates a sharp discontinuity at the boundary. The UNet interprets this boundary as a structural edge, generating persistent incorrect artifacts in subsequent denoising steps.

\paragraph{2. Direct latent blending}
Rather than matching noise levels, this method directly interpolates the clean edit into the noisy control latent using a time-decaying weight $\alpha$ \citep{avrahami2022blended}:
\begin{align*}
  z_t^{\mathrm{new}} &= \bigl(1 - \alpha \cdot m\bigr)\odot z_t^{\mathrm{ctrl}} + \alpha \cdot m \odot \hat{z}_0^{\mathrm{edit}}, \\
  &\quad \text{where } \alpha = 1 - t_{\mathrm{norm}}.
\end{align*}
\textbf{Failure:} Because $\hat{z}_0^{\mathrm{edit}}$ represents a fully denoised state ($t{=}0$) and $z_t^{\mathrm{ctrl}}$ is noisy, their linear combination falls outside the expected distribution of the diffusion process. Consequently, the UNet under-denoises the injected region, yielding flat, over-saturated patches. This degradation is most severe at earlier generation stages (lower $t_{\mathrm{norm}}$), which is precisely when \guardpaint{} interventions are most critical.

\paragraph{3. DDIM inversion}
This strategy attempts to find a compatible noise state by inverting $\hat{z}_0^{\mathrm{edit}}$ forward to timestep $t$ using $d{=}10$ deterministic DDIM steps \citep{song2020denoising} prior to blending:
\begin{align*}
  \hat{z}_t^{\mathrm{inv}} &= \mathrm{DDIMInv}\!\bigl(\hat{z}_0^{\mathrm{edit}},\,t,\,d\bigr), \\
  z_t^{\mathrm{new}} &= (1-\alpha\cdot m)\odot z_t^{\mathrm{ctrl}} + \alpha\cdot m \odot \hat{z}_t^{\mathrm{inv}}, \\
  &\quad \text{where } \alpha = 1-t_{\mathrm{norm}}.
\end{align*}
\textbf{Failure:} DDIM inversion under classifier-free guidance introduces reconstruction errors that scale inversely with $d$ \citep{mokady2023null}. More problematically, the inversion is conditioned on the original adversarial text prompt $c_{\mathrm{text}}$. This conditioning pulls the latent back toward the unsafe concept, actively fighting the safety correction achieved by the tournament. Adjusting $d$ offers no viable compromise: small values produce severe boundary seams, while larger values drastically increase latency without matching the quality of null-text methods.

\paragraph{4. Null-text inversion (selected method).}
We adapt null-text inversion \citep{mokady2023null} to function mid-trajectory. Rather than altering the latents directly, we freeze the UNet $f_\theta$ and the text conditioning $c_{\mathrm{text}}$, and optimize a learnable unconditional embedding $\boldsymbol{\emptyset}^*$. The objective forces the UNet's single-step prediction to align perfectly with the clean edit:
\begin{equation*}
  \boldsymbol{\emptyset}^{*} = \arg\min_{\boldsymbol{\emptyset}} \Bigl\| \hat{z}_0\!\bigl(\hat{z}_0^{\mathrm{edit}},\,t,\,\boldsymbol{\emptyset},\,c_{\mathrm{text}}\bigr) - \hat{z}_0^{\mathrm{edit}} \Bigr\|_2^2
\end{equation*}
where $\hat{z}_0(\cdot)$ denotes the UNet's predicted clean latent under CFG \citep{ho2022classifier}. This objective is minimized for 10 AdamW steps ($\mathrm{lr}{=}0.01$). The latent is then updated via feathered mask blending:
\begin{align*}
  z_t^{\mathrm{new}} &= \bigl(1-\alpha \cdot m\bigr)\odot z_t^{\mathrm{ctrl}} + \alpha \cdot m \odot \hat{z}_0^{\mathrm{edit}}, \\ \qquad
  &\quad \text{where } \alpha = 1-t_{\mathrm{norm}}.
\end{align*}

\textbf{Why null-text inversion succeeds.}
 UNet's prediction becomes self-consistent as a single denoising step from the blended latent reproduces $\hat{z}_0^{\mathrm{edit}}$ inside the mask and preserves the control trajectory outside it. This resolves the independent noise discontinuities of DDPM blending, the state mismatch of direct blending, and the adversarial concept leakage of DDIM inversion. Furthermore, the brief 10-step optimization overhead is incurred only once per audited timestep and is amortized efficiently across all candidates via a per-timestep embedding cache.

\textbf{Flow-matching architectures.}
For flow-matching models (e.g., SD\,3.5 and FLUX.1-dev) where the DDPM forward process does not apply, the noise-matched latent is approximated as $(1-\sigma_t)\hat{z}_0^{\mathrm{edit}} + \sigma_t\boldsymbol{\varepsilon}$ \citep{liu2023flow}, where $\sigma_t$ represents the schedule flow sigma. Adapting native null-text inversion to flow-matching paradigms remains an area for future work.

\clearpage
\newpage

\begin{figure*}[t]
  \centering
  \includegraphics[width=\linewidth]{reinsert-quality.png}
  \caption{
    \textbf{Shows the qualitative difference between different
    reinsertion strategies tried for Stable Diffusion~1.5 and
    SDXL~0.9 Base.}
  }
  \label{fig:rein}
\end{figure*}


%% ────────────────────────────────────────────────────────────────────────────
\subsection{Full Training Configuration}
\label{app:tspo_config}
%% ────────────────────────────────────────────────────────────────────────────

\begin{table}[h]
\centering
\small
\setlength{\tabcolsep}{4pt}
\begin{tabular}{lll}
\toprule
\textbf{Parameter} & \textbf{Value} & \textbf{Notes} \\
\midrule
\multicolumn{3}{l}{\textit{Policy network}} \\
State dim $d_s$   & 257 & $4 \times 64 + 1$ \\
Projection dim    & 64  & Per input \\
Hidden dims       & $(256, 128, 64)$ & LayerNorm + SiLU \\
Log-std clamp     & $[-4.0, +0.5]$ & \\
Weight init       & Orthogonal, gain $0.01$ & \\
\midrule
\multicolumn{3}{l}{\textit{Tournament}} \\
Candidates $N$    & 5   & \\
Audit start       & 70\% of schedule & \\
Audit resolution  & 224 px ($t < 0.65$), 384 px ($t \geq 0.65$) & Coarse/fine \\
$\delta$ init     & 0.05 & Recalibrated every 100 steps \\
$\delta$ min samples & 50 & Before first recalibration \\
\midrule
\multicolumn{3}{l}{\textit{GTP loss}} \\
$\lambda_H^\text{cont}$ & 0.01 & Continuous entropy \\
$\lambda_H^\text{disc}$ & 0.005 & Discrete entropy \\
$\lambda_c$             & 0.005 & Compute penalty \\
$\lambda_\text{div}$    & 0.01 & Diversity \\
$\lambda_\text{listnet}$& 1.0  & ListNet weight in judge \\
$\lambda_\text{PL}$     & 0.5  & Plackett-Luce weight in judge \\
\midrule
\multicolumn{3}{l}{\textit{Optimization}} \\
Policy optimizer  & AdamW, lr $3\!\times\!10^{-4}$ & \\
Judge optimizer   & AdamW, lr $1\!\times\!10^{-4}$ & \\
Mini-batch size   & 4 tournaments (policy), 8 (judge) & \\
Gradient clipping & 1.0 & Both policy and judge \\
Total steps       & 1000 (convergence $\approx$400--600) & \\
\midrule
\multicolumn{3}{l}{\textit{Seam quality}} \\
Ring width        & 8 px & \\
LPIPS net         & VGG (spatial), $\kappa=5.0$ & \\
\midrule
\multicolumn{3}{l}{\textit{Time-varying thresholds}} \\
$\tau_P(t)$       & $0.40 + 0.25\,t$ & \\
$\tau_F(t)$       &$
\tau_F(t_{\text{norm}})=
\begin{cases}
0.30 + 0.30\, t_{\text{norm}}, \\
\quad t_{\text{norm}} < 0.85, \\[8pt]
0.55 - 0.10\,
\left(
\frac{t_{\text{norm}} - 0.85}{0.15}
\right), \\
\quad t_{\text{norm}} \geq 0.85.\notag
\end{cases}
$                 & Non-monotone \\
$\alpha(t)$       & $0.30 + 0.60\,t$ & Blend coefficient \\
\bottomrule
\end{tabular}
\caption{
  Complete Guarded Tournament Policy training and inference configuration.
}
\label{tab:tspo_config}
\end{table}


\end{document}